\documentclass[11pt]{article}
\usepackage[a4paper,margin=1in]{geometry}
\usepackage{amsmath,amssymb,amsfonts,amsthm,mathtools,bm}
\usepackage{booktabs}
\usepackage{array}
\usepackage{longtable}
\usepackage{enumitem}
\usepackage{hyperref}
\usepackage{graphicx}
\usepackage{float}
\usepackage[T1]{fontenc}
\usepackage[utf8]{inputenc}
\hypersetup{colorlinks=true,linkcolor=blue,citecolor=blue,urlcolor=blue}
\linespread{1.08}

\newcommand{\ii}{\mathrm{i}}
\newcommand{\ee}{\mathrm{e}}
\newcommand{\dd}{\mathrm{d}}
\newcommand{\Tr}{\operatorname{Tr}}
\newcommand{\diag}{\operatorname{diag}}
\newcommand{\sgn}{\operatorname{sgn}}
\newcommand{\im}{\operatorname{Im}}
\newcommand{\re}{\operatorname{Re}}
\newcommand{\const}{\mathrm{const}}
\newcommand{\G}{\mathcal{G}}
\newcommand{\K}{\mathcal{K}}
\newcommand{\F}{\mathcal{F}}
\newcommand{\order}[1]{\mathcal{O}\!\left(#1\right)}
\newcommand{\braket}[1]{\left\langle #1 \right\rangle}
\newcommand{\paren}[1]{\left( #1 \right)}
\newcommand{\sqbrak}[1]{\left[ #1 \right]}
\newcommand{\cbrak}[1]{\left\{ #1 \right\}}
\newcommand{\abs}[1]{\left| #1 \right|}
\newcommand{\norm}[1]{\left\| #1 \right\|}

\title{Detailed Theory Notes for Phonon NEGF with Anharmonic Interactions\\\large General \texorpdfstring{$n$}{n}-Phonon Processes, Approximation Hierarchies, and Recovery of the Main Formulas in the Local Literature Set}
\author{negfpy theory workspace}
\date{April 24, 2026}

\begin{document}
\maketitle
\tableofcontents
\newpage

\section*{Purpose and Scope}
These notes are meant to serve two roles at once.
First, they collect a \emph{generalizable} theoretical foundation for phonon nonequilibrium Green's functions (NEGF), written in a way that can later support material-specific implementations with realistic interatomic force constants (IFCs). Second, they record the \emph{benchmark-facing} derivations needed to reproduce and audit the toy-model literature currently stored in \texttt{theory/literature/}.

The emphasis is therefore deliberately two-layered:
\begin{enumerate}[label=\arabic*.]
\item a general derivation of harmonic and anharmonic phonon NEGF, including general $n$-phonon interactions and the practical approximation schemes LO, MF, and SCBA;
\item an explicit recovery of the main formulas used in the local paper set:
\begin{itemize}
\item \texttt{0605028v1.pdf}
\item \texttt{0701164v1.pdf}
\item \texttt{0802.2761v1.pdf}
\item \texttt{0807.3819v1.pdf}
\item \texttt{1303.7317v1.pdf}
\item \texttt{transp.pdf} (journal version of the 2008 review)
\end{itemize}
\end{enumerate}

A recurring theme in these notes is the distinction between
\begin{itemize}
\item \textbf{general theory} that should survive the transition from toy chains to real materials, and
\item \textbf{paper-specific benchmark conventions} that are useful for reproducing a figure but should not be baked into the reusable solver layer.
\end{itemize}

\section{Foundations of Phonon NEGF}

\subsection{Mass-normalized coordinates and the harmonic Hamiltonian}
We begin from physical displacement coordinates $x_i$ and atomic masses $m_i$. The physical kinetic energy is
\begin{equation}
    T = \frac{1}{2}\sum_i m_i \dot x_i^2.
    \label{eq:physical_kinetic_energy}
\end{equation}
Introduce mass-normalized coordinates
\begin{equation}
    u_i = \sqrt{m_i}\,x_i,
    \qquad
    x_i = \frac{u_i}{\sqrt{m_i}}.
    \label{eq:mass_normalized_coordinate_def}
\end{equation}
Differentiate Eq.~\eqref{eq:mass_normalized_coordinate_def} with respect to time:
\begin{equation}
    \dot x_i = \frac{\dot u_i}{\sqrt{m_i}}.
    \label{eq:xdot_from_udot}
\end{equation}
Substitute Eq.~\eqref{eq:xdot_from_udot} into Eq.~\eqref{eq:physical_kinetic_energy}:
\begin{align}
    T
    &= \frac{1}{2}\sum_i m_i \paren{\frac{\dot u_i}{\sqrt{m_i}}}^2 \\
    &= \frac{1}{2}\sum_i m_i \frac{\dot u_i^2}{m_i} \\
    &= \frac{1}{2}\sum_i \dot u_i^2.
    \label{eq:kinetic_energy_mass_normalized}
\end{align}
Thus the masses disappear from the kinetic term.

Now expand the harmonic potential energy around equilibrium in physical coordinates:
\begin{equation}
    V^{(2)} = \frac{1}{2}\sum_{ij}\phi^{(2)}_{ij}x_i x_j.
    \label{eq:physical_harmonic_potential}
\end{equation}
Substituting $x_i = u_i/\sqrt{m_i}$ gives
\begin{align}
    V^{(2)}
    &= \frac{1}{2}\sum_{ij}\phi^{(2)}_{ij}\frac{u_i}{\sqrt{m_i}}\frac{u_j}{\sqrt{m_j}} \\
    &= \frac{1}{2}\sum_{ij}\frac{\phi^{(2)}_{ij}}{\sqrt{m_i m_j}}u_i u_j.
\end{align}
Define the mass-normalized harmonic matrix
\begin{equation}
    K_{ij} = \frac{\phi^{(2)}_{ij}}{\sqrt{m_i m_j}}.
    \label{eq:mass_normalized_harmonic_matrix}
\end{equation}
Then the harmonic Hamiltonian becomes
\begin{equation}
    H^{(2)} = \frac{1}{2}\sum_i \dot u_i^2 + \frac{1}{2}\sum_{ij}u_i K_{ij}u_j.
    \label{eq:harmonic_hamiltonian_mass_normalized}
\end{equation}
In matrix form,
\begin{equation}
    H^{(2)} = \frac{1}{2}\dot u^{\mathsf T}\dot u + \frac{1}{2}u^{\mathsf T}K u.
    \label{eq:harmonic_hamiltonian_matrix_form}
\end{equation}

We now derive the equations of motion from Eq.~\eqref{eq:harmonic_hamiltonian_matrix_form}. The canonical momentum conjugate to $u_i$ is
\begin{equation}
    p_i = \frac{\partial L}{\partial \dot u_i} = \dot u_i,
\end{equation}
so Hamilton's equation gives
\begin{equation}
    \dot p_i = -\frac{\partial H}{\partial u_i} = -\sum_j K_{ij}u_j.
\end{equation}
Since $p_i = \dot u_i$, we get
\begin{equation}
    \ddot u_i = -\sum_j K_{ij}u_j.
    \label{eq:eom_component_harmonic}
\end{equation}
Collecting all components,
\begin{equation}
    \ddot u + Ku = 0.
    \label{eq:eom_matrix_harmonic}
\end{equation}
This is one of the major advantages of mass normalization: the harmonic equations become a clean eigenvalue problem for $K$.

We partition the full system into left lead $L$, center $C$, and right lead $R$. For each region $\alpha \in \{L,C,R\}$, the harmonic Hamiltonian is
\begin{equation}
    H_\alpha = \frac{1}{2}\dot u_\alpha^{\mathsf T}\dot u_\alpha + \frac{1}{2}u_\alpha^{\mathsf T}K^\alpha u_\alpha.
    \label{eq:subsystem_harmonic_hamiltonian}
\end{equation}
The full linear Hamiltonian reads
\begin{equation}
    H_0 = H_L + H_C + H_R + u_L^{\mathsf T}V^{LC}u_C + u_C^{\mathsf T}V^{CR}u_R.
    \label{eq:partitioned_harmonic_hamiltonian}
\end{equation}
In block-matrix notation, the full harmonic matrix is
\begin{equation}
    K =
    \begin{pmatrix}
        K^L & V^{LC} & 0 \\
        V^{CL} & K^C & V^{CR} \\
        0 & V^{RC} & K^R
    \end{pmatrix}.
    \label{eq:blockK}
\end{equation}
The central region is finite, while the leads are infinite or semi-infinite. The entire NEGF reduction to a central problem comes from integrating out the lead blocks of Eq.~\eqref{eq:blockK}.

\subsection{Anharmonic Hamiltonian and convention changes}
The Wang-family papers often write the nonlinear Hamiltonian as
\begin{equation}
    H_n = \frac{1}{3}\sum_{ijk}T_{ijk}u_i u_j u_k + \frac{1}{4}\sum_{ijkl}T_{ijkl}u_i u_j u_k u_l + \cdots.
    \label{eq:wang_Hn}
\end{equation}
A more symmetric field-theory convention is
\begin{equation}
    H_n = \sum_{m\ge 3}\frac{1}{m!}\sum_{i_1\cdots i_m}\Phi^{(m)}_{i_1\cdots i_m}u_{i_1}\cdots u_{i_m}.
    \label{eq:symmetric_Hn}
\end{equation}
To translate between them, compare the cubic terms. Equating the cubic pieces of Eq.~\eqref{eq:wang_Hn} and Eq.~\eqref{eq:symmetric_Hn} gives
\begin{equation}
    \frac{1}{3}T_{ijk}u_i u_j u_k = \frac{1}{3!}\Phi^{(3)}_{ijk}u_i u_j u_k.
\end{equation}
Since $3! = 6$, multiplying both sides by $6$ yields
\begin{equation}
    2T_{ijk}u_i u_j u_k = \Phi^{(3)}_{ijk}u_i u_j u_k,
\end{equation}
and therefore
\begin{equation}
    \Phi^{(3)}_{ijk} = 2T_{ijk}.
    \label{eq:phi3_vs_T3}
\end{equation}
Likewise, for quartic terms we equate
\begin{equation}
    \frac{1}{4}T_{ijkl}u_i u_j u_k u_l = \frac{1}{4!}\Phi^{(4)}_{ijkl}u_i u_j u_k u_l.
\end{equation}
Since $4! = 24$, multiplying both sides by $24$ yields
\begin{equation}
    6T_{ijkl}u_i u_j u_k u_l = \Phi^{(4)}_{ijkl}u_i u_j u_k u_l,
\end{equation}
so
\begin{equation}
    \Phi^{(4)}_{ijkl} = 6T_{ijkl}.
    \label{eq:phi4_vs_T4}
\end{equation}
This is a crucial practical point. Two formulas that look different on paper may in fact describe the same physics once the tensor convention is translated properly.

For future material-based work, the most general Hamiltonian to keep in mind is
\begin{equation}
    H = H_0 + \sum_{m\ge 3}\frac{1}{m!}\sum_{i_1\cdots i_m}\Phi^{(m)}_{i_1\cdots i_m}u_{i_1}\cdots u_{i_m}.
    \label{eq:general_hamiltonian}
\end{equation}
Toy-model tensors are then just sparse special cases of the general IFC tensors $\Phi^{(m)}$.

\subsection{Contour-ordered Green's functions}
The contour-ordered Green's function on the Keldysh contour $C$ is defined by
\begin{equation}
    G_{jk}(\tau,\tau') = -\frac{\ii}{\hbar}\braket{T_C u_j(\tau)u_k(\tau')}.
    \label{eq:contour_green_definition}
\end{equation}
To connect this with real-time transport, we introduce the usual components. If both contour times lie on the forward branch, the time-ordering operator becomes the ordinary real-time ordering operator and one obtains the time-ordered Green's function. From the contour construction we define the lesser and greater functions as
\begin{align}
    G^<_{jk}(t,t') &= -\frac{\ii}{\hbar}\braket{u_k(t')u_j(t)},
    \label{eq:lesser_def} \\
    G^>_{jk}(t,t') &= -\frac{\ii}{\hbar}\braket{u_j(t)u_k(t')}.
    \label{eq:greater_def}
\end{align}
The retarded and advanced functions are then
\begin{align}
    G^r_{jk}(t,t') &= -\frac{\ii}{\hbar}\theta(t-t')\braket{[u_j(t),u_k(t')]},
    \label{eq:retarded_def} \\
    G^a_{jk}(t,t') &= \frac{\ii}{\hbar}\theta(t'-t)\braket{[u_j(t),u_k(t')]}.
    \label{eq:advanced_def}
\end{align}
These are not independent. Start from the commutator identity
\begin{equation}
    [u_j(t),u_k(t')] = u_j(t)u_k(t') - u_k(t')u_j(t).
\end{equation}
Multiply by $-\ii/\hbar$:
\begin{equation}
    -\frac{\ii}{\hbar}\braket{[u_j(t),u_k(t')]} = G^>_{jk}(t,t') - G^<_{jk}(t,t').
\end{equation}
Comparing with Eq.~\eqref{eq:retarded_def} and Eq.~\eqref{eq:advanced_def} yields
\begin{equation}
    G^r - G^a = G^> - G^<.
    \label{eq:spectral_identity}
\end{equation}
In a steady state, all two-time functions depend only on the time difference $t-t'$. We therefore Fourier transform with the convention
\begin{equation}
    G[\omega] = \int_{-\infty}^{+\infty}\dd t\;\ee^{\ii\omega t}G(t),
    \qquad
    G(t) = \int_{-\infty}^{+\infty}\frac{\dd\omega}{2\pi}\;\ee^{-\ii\omega t}G[\omega].
    \label{eq:fourier_convention_GF}
\end{equation}
The sign convention matters because it determines how derivatives transform. Under Eq.~\eqref{eq:fourier_convention_GF}, one has
\begin{equation}
    \frac{\dd}{\dd t} \longrightarrow -\ii\omega,
    \qquad
    \frac{\dd^2}{\dd t^2} \longrightarrow -\omega^2.
    \label{eq:derivative_fourier_rule}
\end{equation}

\subsection{Equation-of-motion derivation of the harmonic retarded Green's function}
We now derive the harmonic retarded Green's function step by step. Start from the definition
\begin{equation}
    G^r_{jk}(t,t') = -\frac{\ii}{\hbar}\theta(t-t')\braket{[u_j(t),u_k(t')]}. 
    \label{eq:retarded_start_eom}
\end{equation}
Differentiate Eq.~\eqref{eq:retarded_start_eom} once with respect to $t$. Use the product rule:
\begin{align}
    \partial_t G^r_{jk}(t,t')
    &= -\frac{\ii}{\hbar}\partial_t\theta(t-t')\braket{[u_j(t),u_k(t')]} 
       -\frac{\ii}{\hbar}\theta(t-t')\partial_t\braket{[u_j(t),u_k(t')]}. 
\end{align}
Since $\partial_t\theta(t-t') = \delta(t-t')$ and only the first operator depends on $t$, this becomes
\begin{align}
    \partial_t G^r_{jk}(t,t')
    &= -\frac{\ii}{\hbar}\delta(t-t')\braket{[u_j(t),u_k(t')]} 
       -\frac{\ii}{\hbar}\theta(t-t')\braket{[\dot u_j(t),u_k(t')]}. 
    \label{eq:first_derivative_retarded}
\end{align}
At equal times, displacement operators commute, so
\begin{equation}
    [u_j(t),u_k(t)] = 0,
\end{equation}
which means the $\delta$-term of Eq.~\eqref{eq:first_derivative_retarded} vanishes. Thus
\begin{equation}
    \partial_t G^r_{jk}(t,t') = -\frac{\ii}{\hbar}\theta(t-t')\braket{[\dot u_j(t),u_k(t')]}. 
    \label{eq:first_derivative_retarded_simplified}
\end{equation}
Differentiate Eq.~\eqref{eq:first_derivative_retarded_simplified} once more:
\begin{align}
    \partial_t^2 G^r_{jk}(t,t')
    &= -\frac{\ii}{\hbar}\delta(t-t')\braket{[\dot u_j(t),u_k(t')]} 
       -\frac{\ii}{\hbar}\theta(t-t')\braket{[\ddot u_j(t),u_k(t')]}. 
    \label{eq:second_derivative_retarded}
\end{align}
Now use the canonical commutator in mass-normalized coordinates. Since the conjugate momentum is $p_j = \dot u_j$, one has
\begin{equation}
    [u_j(t),p_k(t)] = \ii\hbar\delta_{jk}.
\end{equation}
Therefore
\begin{equation}
    [\dot u_j(t),u_k(t)] = [p_j(t),u_k(t)] = -\ii\hbar\delta_{jk}.
    \label{eq:canonical_commutator_udot_u}
\end{equation}
Substitute Eq.~\eqref{eq:canonical_commutator_udot_u} into the first term of Eq.~\eqref{eq:second_derivative_retarded}:
\begin{align}
    -\frac{\ii}{\hbar}\delta(t-t')\braket{[\dot u_j(t),u_k(t')]}
    &= -\frac{\ii}{\hbar}\delta(t-t')\paren{-\ii\hbar\delta_{jk}} \\
    &= -\delta_{jk}\delta(t-t').
\end{align}
Thus Eq.~\eqref{eq:second_derivative_retarded} becomes
\begin{equation}
    \partial_t^2 G^r_{jk}(t,t')
    = -\delta_{jk}\delta(t-t') -\frac{\ii}{\hbar}\theta(t-t')\braket{[\ddot u_j(t),u_k(t')]}. 
    \label{eq:second_derivative_with_delta}
\end{equation}
Use the harmonic equation of motion Eq.~\eqref{eq:eom_component_harmonic}, namely
\begin{equation}
    \ddot u_j(t) = -\sum_l K_{jl}u_l(t).
\end{equation}
Then
\begin{align}
    -\frac{\ii}{\hbar}\theta(t-t')\braket{[\ddot u_j(t),u_k(t')]}
    &= -\frac{\ii}{\hbar}\theta(t-t')\braket{\left[-\sum_l K_{jl}u_l(t),u_k(t')\right]} \\
    &= \sum_l K_{jl}\frac{\ii}{\hbar}\theta(t-t')\braket{[u_l(t),u_k(t')]} \\
    &= -\sum_l K_{jl}G^r_{lk}(t,t').
\end{align}
Hence
\begin{equation}
    \partial_t^2 G^r_{jk}(t,t') = -\delta_{jk}\delta(t-t') - \sum_l K_{jl} G^r_{lk}(t,t').
\end{equation}
Writing this in matrix form,
\begin{equation}
    \sqbrak{\partial_t^2 I + K}G^r(t,t') = -I\delta(t-t').
    \label{eq:EOM_retarded_time}
\end{equation}
Now Fourier transform Eq.~\eqref{eq:EOM_retarded_time}. Using Eq.~\eqref{eq:derivative_fourier_rule},
\begin{equation}
    \partial_t^2 G^r(t,t') \longrightarrow -\omega^2 G^r(\omega).
\end{equation}
The retarded boundary condition is implemented by shifting $\omega \to \omega+\ii\eta$, so the Fourier-space equation is
\begin{equation}
    \sqbrak{-(\omega+\ii\eta)^2I + K}G^r(\omega) = -I.
\end{equation}
Multiply both sides by $-1$:
\begin{equation}
    \sqbrak{(\omega+\ii\eta)^2I - K}G^r(\omega) = I.
\end{equation}
Finally,
\begin{equation}
    G^r(\omega) = \sqbrak{(\omega+\ii\eta)^2I - K}^{-1}.
    \label{eq:retarded_full_harmonic}
\end{equation}
This is the fundamental harmonic retarded Green's function.

\subsection{Eliminating the leads and deriving the central harmonic Green's function}
We now reduce the full harmonic problem to an effective problem in the center. The full resolvent equation is
\begin{equation}
    \sqbrak{(\omega+\ii\eta)^2 I - K}G^r = I.
\end{equation}
Use the block structure of Eq.~\eqref{eq:blockK}. Then
\begin{equation}
    (\omega+\ii\eta)^2I - K =
    \begin{pmatrix}
        (\omega+\ii\eta)^2I - K^L & -V^{LC} & 0 \\
        -V^{CL} & (\omega+\ii\eta)^2I - K^C & -V^{CR} \\
        0 & -V^{RC} & (\omega+\ii\eta)^2I - K^R
    \end{pmatrix}.
    \label{eq:block_resolvent_matrix}
\end{equation}
Write the full Green's function in block form:
\begin{equation}
    G^r =
    \begin{pmatrix}
        G_{LL}^r & G_{LC}^r & G_{LR}^r \\
        G_{CL}^r & G_{CC}^r & G_{CR}^r \\
        G_{RL}^r & G_{RC}^r & G_{RR}^r
    \end{pmatrix}.
\end{equation}
Multiply Eq.~\eqref{eq:block_resolvent_matrix} by this block matrix and focus on the center column. The left and right block equations are
\begin{align}
    \sqbrak{(\omega+\ii\eta)^2I - K^L}G_{LC}^r - V^{LC}G_{CC}^r &= 0,
    \label{eq:left_block_center_column} \\
    \sqbrak{(\omega+\ii\eta)^2I - K^R}G_{RC}^r - V^{RC}G_{CC}^r &= 0.
    \label{eq:right_block_center_column}
\end{align}
Define the isolated lead Green's functions
\begin{equation}
    g_\alpha^r = \sqbrak{(\omega+\ii\eta)^2I - K^\alpha}^{-1},
    \qquad \alpha \in \{L,R\}.
    \label{eq:isolated_lead_retarded}
\end{equation}
Multiply Eq.~\eqref{eq:left_block_center_column} and Eq.~\eqref{eq:right_block_center_column} by $g_L^r$ and $g_R^r$ respectively:
\begin{align}
    G_{LC}^r &= g_L^r V^{LC}G_{CC}^r,
    \label{eq:GLC_in_terms_of_GCC} \\
    G_{RC}^r &= g_R^r V^{RC}G_{CC}^r.
    \label{eq:GRC_in_terms_of_GCC}
\end{align}
Now examine the center block equation:
\begin{equation}
    -V^{CL}G_{LC}^r + \sqbrak{(\omega+\ii\eta)^2I - K^C}G_{CC}^r - V^{CR}G_{RC}^r = I.
    \label{eq:center_block_equation_raw}
\end{equation}
Insert Eq.~\eqref{eq:GLC_in_terms_of_GCC} and Eq.~\eqref{eq:GRC_in_terms_of_GCC} into Eq.~\eqref{eq:center_block_equation_raw}:
\begin{align}
    &-V^{CL}\paren{g_L^r V^{LC}G_{CC}^r} + \sqbrak{(\omega+\ii\eta)^2I - K^C}G_{CC}^r - V^{CR}\paren{g_R^r V^{RC}G_{CC}^r} = I.
\end{align}
Factor out $G_{CC}^r$ on the right:
\begin{equation}
    \sqbrak{(\omega+\ii\eta)^2I - K^C - V^{CL}g_L^rV^{LC} - V^{CR}g_R^rV^{RC}}G_{CC}^r = I.
    \label{eq:center_block_after_elimination}
\end{equation}
Define the lead self-energies
\begin{equation}
    \Sigma_\alpha^r(\omega) = V^{C\alpha}g_\alpha^r(\omega)V^{\alpha C},
    \label{eq:lead_self_energy_definition}
\end{equation}
where $V^{CL} \equiv V^{C L}$ and $V^{LC} \equiv V^{L C}$, etc. Then Eq.~\eqref{eq:center_block_after_elimination} becomes
\begin{equation}
    G_0^r(\omega) \equiv G_{CC}^r(\omega)
    = \sqbrak{(\omega+\ii\eta)^2I - K^C - \Sigma_L^r(\omega) - \Sigma_R^r(\omega)}^{-1}.
    \label{eq:G0r}
\end{equation}
This is the exact harmonic center Green's function after the leads have been integrated out.

\subsection{Lesser Green's function and Keldysh relation}
In the harmonic problem, nonequilibrium enters only through the thermal occupations of the leads. The contour Dyson equation for the center is
\begin{equation}
    G_0 = g_C + g_C \otimes \Sigma \otimes G_0,
    \label{eq:contour_dyson_harmonic_center}
\end{equation}
where $\Sigma = \Sigma_L + \Sigma_R$ and $\otimes$ denotes contour convolution.

Apply Langreth rules to Eq.~\eqref{eq:contour_dyson_harmonic_center}. The lesser component of a product obeys
\begin{equation}
    (AB)^< = A^r B^< + A^< B^a.
    \label{eq:langreth_product_rule}
\end{equation}
Using this repeatedly on the convolution form of Eq.~\eqref{eq:contour_dyson_harmonic_center} and assuming the isolated center carries no initial nonequilibrium population, one arrives at
\begin{equation}
    G_0^< = G_0^r \Sigma^< G_0^a.
    \label{eq:G0lesser}
\end{equation}
To express $\Sigma_\alpha^<$ in terms of physical lead quantities, define the level-width matrix
\begin{equation}
    \Gamma_\alpha(\omega) = \ii\paren{\Sigma_\alpha^r - \Sigma_\alpha^a}.
    \label{eq:Gamma_definition}
\end{equation}
For a lead in local equilibrium at temperature $T_\alpha$, the lesser self-energy is
\begin{equation}
    \Sigma_\alpha^<(\omega) = f_\alpha(\omega)\paren{\Sigma_\alpha^a - \Sigma_\alpha^r}.
\end{equation}
Using Eq.~\eqref{eq:Gamma_definition},
\begin{equation}
    \Sigma_\alpha^a - \Sigma_\alpha^r = -\paren{\Sigma_\alpha^r - \Sigma_\alpha^a} = -\frac{1}{\ii}\Gamma_\alpha = -\ii\Gamma_\alpha,
\end{equation}
so
\begin{equation}
    \Sigma_\alpha^<(\omega) = -\ii f_\alpha(\omega)\Gamma_\alpha(\omega).
    \label{eq:lead_lesser_self_energy_bose}
\end{equation}
The Bose distribution is
\begin{equation}
    f_\alpha(\omega) = \frac{1}{\exp\sqbrak{\hbar\omega/(k_B T_\alpha)} - 1}.
    \label{eq:bose_distribution}
\end{equation}
Combining Eq.~\eqref{eq:G0lesser} and Eq.~\eqref{eq:lead_lesser_self_energy_bose}, we obtain
\begin{equation}
    G_0^< = G_0^r\paren{\Sigma_L^< + \Sigma_R^<}G_0^a.
    \label{eq:harmonic_lesser_final}
\end{equation}
This is the harmonic Keldysh relation used in both ballistic and inelastic calculations.

\subsection{Dyson and Keldysh equations with anharmonicity}
Once anharmonicity is added, the lead self-energies are supplemented by a nonlinear self-energy $\Sigma_n$. The exact interacting central-region Dyson equation is
\begin{equation}
    G^r = G_0^r + G_0^r\Sigma_n^r G^r.
    \label{eq:dyson_interacting_integral_form}
\end{equation}
Bring the second term to the left:
\begin{equation}
    G^r - G_0^r\Sigma_n^r G^r = G_0^r.
\end{equation}
Factor out $G^r$ on the right:
\begin{equation}
    \paren{I - G_0^r\Sigma_n^r}G^r = G_0^r.
\end{equation}
Multiply by $(G_0^r)^{-1}$ from the left:
\begin{equation}
    \paren{(G_0^r)^{-1} - \Sigma_n^r}G^r = I.
\end{equation}
Hence
\begin{equation}
    G^r = \sqbrak{(G_0^r)^{-1} - \Sigma_n^r}^{-1}.
    \label{eq:Dyson_r}
\end{equation}
This is Dyson's equation in algebraic form.

For the lesser component, the same Langreth machinery yields
\begin{equation}
    G^< = G^r\paren{\Sigma_L^< + \Sigma_R^< + \Sigma_n^<}G^a.
    \label{eq:Keldysh_less}
\end{equation}
Equation~\eqref{eq:Dyson_r} and Eq.~\eqref{eq:Keldysh_less} are exact in structure. Every approximation scheme later in the notes should be understood as a prescription for approximating $\Sigma_n^r$ and $\Sigma_n^<$.

\subsection{Heat-current formula from the lead energy derivative}
We now derive the current formula in detail. Define the current leaving the left lead by
\begin{equation}
    I_L = -\braket{\dot H_L}.
    \label{eq:current_def_start}
\end{equation}
Use the Heisenberg equation of motion,
\begin{equation}
    \dot H_L = \frac{\ii}{\hbar}[H,H_L].
\end{equation}
Since $H_L$ commutes with itself and with all center or right-lead terms except the coupling $u_L^{\mathsf T}V^{LC}u_C$, only the coupling contributes. Differentiating explicitly from the classical-looking form gives the same result more directly:
\begin{equation}
    H_L = \frac{1}{2}\dot u_L^{\mathsf T}\dot u_L + \frac{1}{2}u_L^{\mathsf T}K^L u_L.
\end{equation}
Then
\begin{equation}
    \dot H_L = \dot u_L^{\mathsf T}\ddot u_L + \dot u_L^{\mathsf T}K^L u_L.
\end{equation}
Use the left-lead equation of motion in the full coupled system,
\begin{equation}
    \ddot u_L = -K^L u_L - V^{LC}u_C.
\end{equation}
Substitute this into $\dot H_L$:
\begin{align}
    \dot H_L
    &= \dot u_L^{\mathsf T}\paren{-K^L u_L - V^{LC}u_C} + \dot u_L^{\mathsf T}K^L u_L \\
    &= -\dot u_L^{\mathsf T}V^{LC}u_C.
\end{align}
Therefore
\begin{equation}
    I_L = -\braket{\dot H_L} = \braket{\dot u_L^{\mathsf T}V^{LC}u_C}.
    \label{eq:current_coordinate_velocity_form}
\end{equation}
Now express Eq.~\eqref{eq:current_coordinate_velocity_form} in terms of a mixed lesser Green's function. Define
\begin{equation}
    G_{CL}^<(t,t') = -\frac{\ii}{\hbar}\braket{u_L(t')u_C^{\mathsf T}(t)}.
    \label{eq:mixed_lesser_definition}
\end{equation}
Set $t=t'$ after differentiating with respect to the lead time. Since in frequency space $\dot u_L[\omega] = -\ii\omega u_L[\omega]$, Eq.~\eqref{eq:current_coordinate_velocity_form} becomes
\begin{equation}
    I_L = -\int_{-\infty}^{+\infty}\frac{\dd\omega}{2\pi}\,\hbar\omega\,\Tr\sqbrak{V^{LC}G_{CL}^<(\omega)}.
    \label{eq:current_mixed}
\end{equation}
We now eliminate $G_{CL}^<$ in favor of center quantities. On the contour, the mixed Green's function obeys
\begin{equation}
    G_{CL} = G V^{CL} g_L,
\end{equation}
where $G$ is the full center Green's function and $g_L$ is the isolated lead Green's function. Applying Langreth rules to the product gives
\begin{equation}
    G_{CL}^< = G^r V^{CL}g_L^< + G^<V^{CL}g_L^a.
    \label{eq:GCLlesser_langreth}
\end{equation}
Insert Eq.~\eqref{eq:GCLlesser_langreth} into Eq.~\eqref{eq:current_mixed}:
\begin{align}
    I_L
    &= -\int_{-\infty}^{+\infty}\frac{\dd\omega}{2\pi}\,\hbar\omega\,
    \Tr\sqbrak{V^{LC}G^r V^{CL}g_L^< + V^{LC}G^<V^{CL}g_L^a}.
\end{align}
Use cyclicity of the trace and the definitions
\begin{equation}
    \Sigma_L^< = V^{CL}g_L^<V^{LC},
    \qquad
    \Sigma_L^a = V^{CL}g_L^aV^{LC}.
\end{equation}
Then
\begin{equation}
    I_L = -\int_{-\infty}^{+\infty}\frac{\dd\omega}{2\pi}\,\hbar\omega\,
    \Tr\sqbrak{G^r\Sigma_L^< + G^<\Sigma_L^a}.
    \label{eq:current_left}
\end{equation}
This is the exact interacting phonon current formula used throughout the Wang-family papers.

\subsection{Ballistic Landauer--Caroli formula}
In the purely harmonic case, $\Sigma_n = 0$, so Eq.~\eqref{eq:Keldysh_less} reduces to Eq.~\eqref{eq:harmonic_lesser_final}:
\begin{equation}
    G^< = G_0^r(\Sigma_L^< + \Sigma_R^<)G_0^a.
    \label{eq:ballistic_lesser_repeat}
\end{equation}
Insert Eq.~\eqref{eq:ballistic_lesser_repeat} into Eq.~\eqref{eq:current_left}:
\begin{align}
    I_L
    &= -\int\frac{\dd\omega}{2\pi}\,\hbar\omega\,
    \Tr\sqbrak{G_0^r\Sigma_L^< + G_0^r(\Sigma_L^< + \Sigma_R^<)G_0^a\Sigma_L^a}.
    \label{eq:current_after_ballistic_substitution}
\end{align}
Using standard identities among retarded, advanced, lesser, and broadening matrices, one can rearrange Eq.~\eqref{eq:current_after_ballistic_substitution} into the two-terminal form
\begin{equation}
    I = \int_0^{+\infty}\frac{\dd\omega}{2\pi}\,\hbar\omega\,\mathcal{T}(\omega)\sqbrak{f_L(\omega)-f_R(\omega)},
    \label{eq:landauer_phonon}
\end{equation}
where the transmission is
\begin{equation}
    \mathcal{T}(\omega) = \Tr\sqbrak{G_0^r(\omega)\Gamma_L(\omega)G_0^a(\omega)\Gamma_R(\omega)}.
    \label{eq:caroli}
\end{equation}
Equation~\eqref{eq:caroli} is the Caroli formula. It is the exact ballistic limit of the Green's-function formalism.

\subsection{Thermal conductance}
For a small temperature bias,
\begin{equation}
    T_L = T + \frac{\Delta T}{2},
    \qquad
    T_R = T - \frac{\Delta T}{2}.
    \label{eq:small_bias_temperature_split}
\end{equation}
The thermal conductance is defined by the linear-response slope
\begin{equation}
    \kappa(T) = \lim_{\Delta T \to 0}\frac{I\paren{T+\Delta T/2,\,T-\Delta T/2}}{\Delta T}.
    \label{eq:kappa_definition_general}
\end{equation}
Substitute the ballistic current Eq.~\eqref{eq:landauer_phonon} into Eq.~\eqref{eq:kappa_definition_general}:
\begin{equation}
    \kappa(T) = \lim_{\Delta T\to 0}\int_0^{+\infty}\frac{\dd\omega}{2\pi}\,\hbar\omega\,\mathcal T(\omega)
    \frac{f\paren{\omega,T+\Delta T/2} - f\paren{\omega,T-\Delta T/2}}{\Delta T}.
\end{equation}
Take the limit inside the integral. The bracketed quotient becomes a derivative:
\begin{equation}
    \lim_{\Delta T\to 0}
    \frac{f\paren{\omega,T+\Delta T/2} - f\paren{\omega,T-\Delta T/2}}{\Delta T}
    = \frac{\partial f}{\partial T}.
\end{equation}
Hence
\begin{equation}
    \kappa_{\rm bal}(T) = \int_0^{+\infty}\frac{\dd\omega}{2\pi}\,\hbar\omega\,\mathcal T(\omega)\frac{\partial f}{\partial T}.
    \label{eq:kappa_ballistic}
\end{equation}
This formula is the linear-response conductance benchmark that every inelastic extension should recover when the nonlinear self-energy is turned off.

\section{General \texorpdfstring{$n$}{n}-Phonon Processes, and the Special Cases of Three- and Four-Phonon Interactions}

\subsection{General interaction vertices}
We begin from the Born--Oppenheimer potential energy surface expanded around the chosen equilibrium reference configuration. In physical displacement coordinates $x_i$, the Taylor series is
\begin{equation}
    V(\{x\})
    =
    V_0
    + \frac{1}{2!}\sum_{ij}\phi^{(2)}_{ij}x_i x_j
    + \frac{1}{3!}\sum_{ijk}\phi^{(3)}_{ijk}x_i x_j x_k
    + \frac{1}{4!}\sum_{ijkl}\phi^{(4)}_{ijkl}x_i x_j x_k x_l
    + \cdots.
    \label{eq:pes_taylor_physical}
\end{equation}
The coefficient $\phi^{(m)}_{i_1\cdots i_m}$ is the $m$th derivative of the potential with respect to physical coordinates, evaluated at the reference equilibrium point. Because mixed partial derivatives commute, the tensor is invariant under any permutation of its indices,
\begin{equation}
    \phi^{(m)}_{i_1\cdots i_m} = \phi^{(m)}_{i_{\pi(1)}\cdots i_{\pi(m)}}
    \qquad \text{for any permutation } \pi.
    \label{eq:physical_ifc_permutation_symmetry}
\end{equation}

Now pass to mass-normalized coordinates,
\begin{equation}
    u_i = \sqrt{m_i}x_i,
    \qquad
    x_i = \frac{u_i}{\sqrt{m_i}}.
    \label{eq:mass_norm_for_interactions}
\end{equation}
Substitute Eq.~\eqref{eq:mass_norm_for_interactions} into the $m$th-order term of Eq.~\eqref{eq:pes_taylor_physical}:
\begin{align}
    V^{(m)}
    &= \frac{1}{m!}\sum_{i_1\cdots i_m}\phi^{(m)}_{i_1\cdots i_m}
       x_{i_1}\cdots x_{i_m} \\
    &= \frac{1}{m!}\sum_{i_1\cdots i_m}\phi^{(m)}_{i_1\cdots i_m}
       \frac{u_{i_1}}{\sqrt{m_{i_1}}}\cdots \frac{u_{i_m}}{\sqrt{m_{i_m}}} \\
    &= \frac{1}{m!}\sum_{i_1\cdots i_m}
       \frac{\phi^{(m)}_{i_1\cdots i_m}}{\sqrt{m_{i_1}\cdots m_{i_m}}}
       u_{i_1}\cdots u_{i_m}.
\end{align}
This motivates the mass-normalized anharmonic tensor
\begin{equation}
    \Phi^{(m)}_{i_1\cdots i_m}
    =
    \frac{\phi^{(m)}_{i_1\cdots i_m}}{\sqrt{m_{i_1}\cdots m_{i_m}}}.
    \label{eq:mass_norm_ifc_general_again}
\end{equation}
With this definition the full interaction Hamiltonian is
\begin{equation}
    H_{\rm int}
    =
    \sum_{m\ge 3}
    \frac{1}{m!}
    \sum_{i_1\cdots i_m}
    \Phi^{(m)}_{i_1\cdots i_m}
    u_{i_1}\cdots u_{i_m}.
    \label{eq:Hint_tensor_general}
\end{equation}
Equation~\eqref{eq:Hint_tensor_general} is the most general tensor-level statement of the phonon interaction problem. In a realistic material workflow, this is the object imported from third-order, fourth-order, or higher-order IFC calculations.

When the interaction Hamiltonian is inserted into the contour-ordered evolution operator, each vertex is local in contour time because the IFCs are static numbers. Thus an $m$-leg vertex contributes the contour kernel
\begin{equation}
    \Phi^{(m)}_{i_1\cdots i_m}(\tau_1,\ldots,\tau_m)
    =
    \Phi^{(m)}_{i_1\cdots i_m}
    \delta_C(\tau_1-\tau_2)\cdots \delta_C(\tau_1-\tau_m).
    \label{eq:local_contour_vertex_again}
\end{equation}
This formula is worth emphasizing. It means that all time nonlocality in the self-energy comes from the Green's functions, not from the bare interaction vertex.

\subsection{Normal-mode representation and general scattering channels}
To interpret Eq.~\eqref{eq:Hint_tensor_general} physically, we now transform to the harmonic normal-mode basis. Let the harmonic eigenvalue problem be
\begin{equation}
    \sum_j K_{ij}e_{j\lambda} = \omega_\lambda^2 e_{i\lambda}.
    \label{eq:harmonic_mode_eigenproblem}
\end{equation}
The eigenvectors are chosen orthonormal,
\begin{equation}
    \sum_i e_{i\lambda}e_{i\lambda'} = \delta_{\lambda\lambda'},
    \qquad
    \sum_\lambda e_{i\lambda}e_{j\lambda} = \delta_{ij}.
    \label{eq:harmonic_mode_orthogonality}
\end{equation}
Expand the displacement operators in that basis:
\begin{equation}
    u_i = \sum_\lambda e_{i\lambda}Q_\lambda.
    \label{eq:u_mode_expand_again}
\end{equation}
The normal coordinates are quantized as
\begin{equation}
    Q_\lambda = \sqrt{\frac{\hbar}{2\omega_\lambda}}\paren{a_\lambda + a_\lambda^\dagger}.
    \label{eq:normal_coordinate_quantization_again}
\end{equation}
Now substitute Eq.~\eqref{eq:u_mode_expand_again} into the $m$th-order interaction term:
\begin{align}
    H^{(m)}
    &=
    \frac{1}{m!}\sum_{i_1\cdots i_m}\Phi^{(m)}_{i_1\cdots i_m}
    \paren{\sum_{\lambda_1}e_{i_1\lambda_1}Q_{\lambda_1}}
    \cdots
    \paren{\sum_{\lambda_m}e_{i_m\lambda_m}Q_{\lambda_m}} \\
    &=
    \frac{1}{m!}
    \sum_{\lambda_1\cdots \lambda_m}
    \sum_{i_1\cdots i_m}
    \Phi^{(m)}_{i_1\cdots i_m}
    e_{i_1\lambda_1}\cdots e_{i_m\lambda_m}
    Q_{\lambda_1}\cdots Q_{\lambda_m}.
\end{align}
Define the mode-space interaction tensor
\begin{equation}
    \widetilde\Phi^{(m)}_{\lambda_1\cdots \lambda_m}
    =
    \sum_{i_1\cdots i_m}
    \Phi^{(m)}_{i_1\cdots i_m}
    e_{i_1\lambda_1}\cdots e_{i_m\lambda_m}.
    \label{eq:mode_ifc_not_yet_quantized}
\end{equation}
Then
\begin{equation}
    H^{(m)} = \frac{1}{m!}\sum_{\lambda_1\cdots \lambda_m}\widetilde\Phi^{(m)}_{\lambda_1\cdots \lambda_m}Q_{\lambda_1}\cdots Q_{\lambda_m}.
    \label{eq:mode_hamiltonian_before_aadag}
\end{equation}
Now insert Eq.~\eqref{eq:normal_coordinate_quantization_again} for each normal coordinate. Since every $Q_{\lambda_p}$ contributes a factor $\sqrt{\hbar/(2\omega_{\lambda_p})}$, we get
\begin{align}
    H^{(m)}
    &= \frac{1}{m!}\sum_{\lambda_1\cdots \lambda_m}
       \widetilde\Phi^{(m)}_{\lambda_1\cdots \lambda_m}
       \prod_{p=1}^m \sqrt{\frac{\hbar}{2\omega_{\lambda_p}}}
       \prod_{p=1}^m \paren{a_{\lambda_p}+a_{\lambda_p}^\dagger} \\
    &= \frac{1}{m!}\sum_{\lambda_1\cdots \lambda_m}
       V^{(m)}_{\lambda_1\cdots \lambda_m}
       \prod_{p=1}^m \paren{a_{\lambda_p}+a_{\lambda_p}^\dagger},
    \label{eq:mode_hamiltonian_quantized}
\end{align}
where
\begin{equation}
    V^{(m)}_{\lambda_1\cdots \lambda_m}
    =
    \sum_{i_1\cdots i_m}
    \Phi^{(m)}_{i_1\cdots i_m}
    e_{i_1\lambda_1}\cdots e_{i_m\lambda_m}
    \prod_{p=1}^m \sqrt{\frac{\hbar}{2\omega_{\lambda_p}}}.
    \label{eq:mode_vertex_general}
\end{equation}
Equation~\eqref{eq:mode_vertex_general} is the general mode-space vertex that later appears in lifetime formulas and scattering rates.

To understand the scattering channels, expand the product
\begin{equation}
    \prod_{p=1}^m \paren{a_{\lambda_p}+a_{\lambda_p}^\dagger}.
\end{equation}
Each factor contributes either an annihilation operator or a creation operator, so there are $2^m$ operator monomials. In the interaction picture,
\begin{equation}
    a_\lambda(t) = a_\lambda\ee^{-\ii\omega_\lambda t},
    \qquad
    a_\lambda^\dagger(t) = a_\lambda^\dagger\ee^{+\ii\omega_\lambda t}.
    \label{eq:a_time_evolution_modes}
\end{equation}
Therefore each operator monomial carries a time phase of the form
\begin{equation}
    \exp\sqbrak{-\ii\paren{s_1\omega_{\lambda_1} + \cdots + s_m\omega_{\lambda_m}}t},
    \qquad s_p = \pm 1,
\end{equation}
where $s_p = +1$ corresponds to an annihilation operator and $s_p=-1$ to a creation operator under the sign convention of the phase. The stationary or resonant channels are exactly those for which the phase is time independent, namely
\begin{equation}
    s_1\omega_{\lambda_1} + \cdots + s_m\omega_{\lambda_m} = 0.
    \label{eq:general_mode_energy_conservation}
\end{equation}
For $m=3$, Eq.~\eqref{eq:general_mode_energy_conservation} contains channels such as
\begin{equation}
    \omega_1 - \omega_2 - \omega_3 = 0,
    \qquad
    \omega_1 + \omega_2 - \omega_3 = 0.
\end{equation}
For $m=4$, it contains channels such as
\begin{equation}
    \omega_1 + \omega_2 - \omega_3 - \omega_4 = 0.
\end{equation}
This is the precise meaning of a general $n$-phonon process: it is the collection of all channels encoded by one $n$-leg interaction vertex after the harmonic basis is quantized.

\subsection{General structure of the self-energy}
We now derive the general diagrammatic structure of the self-energy from the contour interaction expansion. Start from the definition of the full contour Green's function,
\begin{equation}
    G_{12}(\tau,\tau')
    =
    -\frac{\ii}{\hbar}
    \frac{\braket{T_C u_1(\tau)u_2(\tau') S_C}}{\braket{S_C}},
    \qquad
    S_C = \exp\sqbrak{-\frac{\ii}{\hbar}\int_C \dd\tau_1 H_{\rm int}(\tau_1)}.
    \label{eq:full_contour_green_with_SC}
\end{equation}
Expand $S_C$ in powers of the interaction,
\begin{equation}
    S_C = 1 - \frac{\ii}{\hbar}\int_C \dd\tau_1 H_{\rm int}(\tau_1)
    + \frac{1}{2!}\paren{-\frac{\ii}{\hbar}}^2\int_C \dd\tau_1\dd\tau_2 H_{\rm int}(\tau_1)H_{\rm int}(\tau_2) + \cdots.
    \label{eq:SC_expansion_general}
\end{equation}
Insert Eq.~\eqref{eq:SC_expansion_general} into Eq.~\eqref{eq:full_contour_green_with_SC}. The result is a series of contour-ordered expectation values evaluated with respect to the harmonic reference system. Wick's theorem reduces each term to sums of products of harmonic propagators. The self-energy is defined as the sum of one-particle-irreducible two-point contributions generated by those contractions.

Rather than jumping directly to the final answer, it is useful to count legs carefully. Suppose a diagram contains $V_m$ vertices of order $m$, and no other vertex types. Then the total number of interaction legs is
\begin{equation}
    N_{\rm legs} = mV_m.
\end{equation}
A two-point self-energy diagram has two external legs, so the number of remaining internal legs is
\begin{equation}
    N_{\rm internal\,legs} = mV_m - 2.
\end{equation}
Every propagator consumes two internal legs. Therefore the number of internal Green's-function lines is
\begin{equation}
    N_G = \frac{mV_m-2}{2}.
    \label{eq:internal_line_count_general}
\end{equation}
Equation~\eqref{eq:internal_line_count_general} immediately gives the two most important cases in phonon NEGF:
\begin{align}
    m=3,\;V_3=2 &\Rightarrow N_G = 2, \\
    m=4,\;V_4=1 &\Rightarrow N_G = 1.
\end{align}
The first line is the leading dynamic cubic self-energy. The second line is the leading static quartic Hartree self-energy.

We now derive the index structure of the static contribution. Take one $m$-leg vertex and keep two indices external, say $1$ and $2$. The remaining $m-2$ operators are contracted into equal-time expectation values. Therefore the resulting two-point object has the structure
\begin{equation}
    \Sigma^{\rm static}_{12}(\tau,\tau')
    \propto
    \delta_C(\tau-\tau')
    \sum_{3\cdots m}
    \Phi^{(m)}_{12 3\cdots m}
    \braket{T_C u_3(\tau)\cdots u_m(\tau)}.
    \label{eq:static_self_energy_general_detail}
\end{equation}
The contour delta function appears because all operators originate from one local-in-time interaction vertex.

For the leading dynamic contribution from two vertices, distribute one external leg on the first vertex and one on the second vertex, then contract all remaining legs pairwise. The result has the general contour structure
\begin{equation}
    \Sigma^{\rm dyn}_{12}(\tau,\tau')
    \propto
    \sum \Phi_{1\cdots} G(\tau,\tau')\cdots G(\tau,\tau')\Phi_{2\cdots},
    \label{eq:dynamic_self_energy_general_detail}
\end{equation}
with the exact number of Green's functions fixed by Eq.~\eqref{eq:internal_line_count_general}. This is the basic origin of the frequency-dependent scattering self-energy.

\subsection{Three-phonon interaction}
For cubic interactions,
\begin{equation}
    H^{(3)} = \frac{1}{3!}\sum_{ijk}\Phi^{(3)}_{ijk}u_i u_j u_k.
    \label{eq:cubic_hamiltonian_general_again}
\end{equation}
We now derive why the lowest nontrivial dynamic self-energy is second order in the cubic vertex. Expand the contour Green's function to first order in Eq.~\eqref{eq:cubic_hamiltonian_general_again}. The first-order correction is
\begin{equation}
    \delta G^{(1)}_{12}
    =
    -\frac{\ii}{\hbar}
    \frac{1}{3!}
    \sum_{abc}\Phi^{(3)}_{abc}
    \int_C \dd\tau_1\,
    \braket{T_C u_1 u_2 u_a u_b u_c}_0.
    \label{eq:first_order_cubic_green_correction}
\end{equation}
The expectation value in Eq.~\eqref{eq:first_order_cubic_green_correction} contains five harmonic fields. For a Gaussian harmonic reference state with vanishing odd moments, every expectation value with an odd number of fields vanishes. Therefore
\begin{equation}
    \delta G^{(1)}_{12} = 0
    \qquad
    \text{for the dynamic cubic correction in the centered harmonic reference state.}
\end{equation}
This is the formal reason the leading cubic dynamic self-energy does not appear at first order.

Now expand to second order in the cubic interaction:
\begin{align}
    \delta G^{(2)}_{12}
    &=
    \frac{1}{2!}\paren{-\frac{\ii}{\hbar}}^2
    \int_C \dd\tau_3\dd\tau_4\,
    \braket{T_C u_1 u_2 H^{(3)}(\tau_3)H^{(3)}(\tau_4)}_0 \\
    &=
    \frac{1}{2!}\paren{-\frac{\ii}{\hbar}}^2
    \frac{1}{(3!)^2}
    \sum_{abc}\sum_{def}
    \Phi^{(3)}_{abc}\Phi^{(3)}_{def}
    \int_C \dd\tau_3\dd\tau_4\,
    \braket{T_C u_1 u_2 u_a u_b u_c u_d u_e u_f}_0.
    \label{eq:second_order_cubic_green_correction}
\end{align}
Now Wick-contract the eight fields. The one-particle-irreducible two-point contribution is the one in which one external operator is attached to the first vertex, the other external operator is attached to the second vertex, and the remaining four internal fields are paired across the two vertices. This gives two internal propagators. After collecting the contractions, the self-energy structure is
\begin{equation}
    \Sigma^{(3)}_{12}(\tau,\tau')
    =
    \frac{\ii\hbar}{2}
    \sum_{3456}
    \Phi^{(3)}_{134}
    G_{35}(\tau,\tau')
    G_{46}(\tau,\tau')
    \Phi^{(3)}_{256}
    + \text{tadpole-type terms},
    \label{eq:cubic_contour_self_energy_detail}
\end{equation}
up to convention-dependent prefactors. The exact coefficient changes with the normalization of the interaction tensor, but the index and propagator structure is fixed. Equation~\eqref{eq:cubic_contour_self_energy_detail} is the general contour expression behind cubic lowest-order scattering.

To connect Eq.~\eqref{eq:cubic_contour_self_energy_detail} with a quasiparticle rate, project it to harmonic mode space and then to the retarded component. Near an isolated mode pole, the dressed propagator denominator is
\begin{equation}
    D_\lambda(\omega) = \omega^2 - \omega_\lambda^2 - \Sigma^{r,(3)}_\lambda(\omega).
    \label{eq:mode_denominator_with_self_energy}
\end{equation}
Suppose the interaction is weak enough that the pole is near $\omega = \omega_\lambda$. Write the complex pole as
\begin{equation}
    \omega = \widetilde\omega_\lambda - \ii\Gamma_\lambda,
    \qquad
    \widetilde\omega_\lambda = \omega_\lambda + \Delta\omega_\lambda.
    \label{eq:complex_mode_pole_definition}
\end{equation}
Insert Eq.~\eqref{eq:complex_mode_pole_definition} into Eq.~\eqref{eq:mode_denominator_with_self_energy} and expand to first order in $\Delta\omega_\lambda$ and $\Gamma_\lambda$:
\begin{align}
    \omega^2
    &= \paren{\omega_\lambda + \Delta\omega_\lambda - \ii\Gamma_\lambda}^2 \\
    &\approx \omega_\lambda^2 + 2\omega_\lambda\Delta\omega_\lambda - 2\ii\omega_\lambda\Gamma_\lambda.
\end{align}
Therefore the pole condition $D_\lambda(\omega)=0$ becomes
\begin{equation}
    2\omega_\lambda\Delta\omega_\lambda - 2\ii\omega_\lambda\Gamma_\lambda - \Sigma^{r,(3)}_\lambda(\omega_\lambda) \approx 0.
\end{equation}
Separating real and imaginary parts gives
\begin{align}
    \Delta\omega_\lambda &\approx \frac{\re\Sigma^{r,(3)}_\lambda(\omega_\lambda)}{2\omega_\lambda}, \\
    \Gamma_\lambda &\approx -\frac{\im\Sigma^{r,(3)}_\lambda(\omega_\lambda)}{2\omega_\lambda}.
    \label{eq:cubic_shift_and_linewidth_again}
\end{align}
Finally,
\begin{equation}
    \tau_\lambda^{-1} = 2\Gamma_\lambda.
    \label{eq:lifetime_from_linewidth_again}
\end{equation}
This is the mode-lifetime formula used in Xu \emph{et al.} 2008.

We now show how the Wang 2007 onsite cubic model emerges as a special case of the general expression. In that model,
\begin{equation}
    H_n^{\rm onsite} = \frac{1}{3}\sum_j t_j u_j^3.
    \label{eq:onsite_cubic_special_hamiltonian}
\end{equation}
Comparing Eq.~\eqref{eq:onsite_cubic_special_hamiltonian} with Eq.~\eqref{eq:cubic_hamiltonian_general_again}, we see that the only nonzero cubic tensor elements are local in site index. Therefore every tensor contraction in Eq.~\eqref{eq:cubic_contour_self_energy_detail} collapses onto one site, and the general expression reduces to
\begin{equation}
    \Sigma_{n,jl}^{\sigma\sigma'}(t) = 2\ii t^2 \sqbrak{G_{0,jl}^{\sigma\sigma'}(t)}^2,
    \label{eq:wang77_recovered_again}
\end{equation}
which is precisely Wang 2007 Eq.~(77). This is why that toy model is such a useful benchmark: it converts the full cubic tensor structure into a single-site scalar interaction while preserving the genuine cubic two-vertex physics.

\subsection{Four-phonon interaction}
For quartic interactions,
\begin{equation}
    H^{(4)} = \frac{1}{4!}\sum_{ijkl}\Phi^{(4)}_{ijkl}u_i u_j u_k u_l.
    \label{eq:quartic_hamiltonian_general_again}
\end{equation}
The reason quartic interactions are naturally associated with mean-field closures is that a nonzero two-point correction already appears at first order. Expand the contour Green's function to first order in Eq.~\eqref{eq:quartic_hamiltonian_general_again}:
\begin{equation}
    \delta G^{(1)}_{12}
    =
    -\frac{\ii}{\hbar}\frac{1}{4!}
    \sum_{abcd}\Phi^{(4)}_{abcd}
    \int_C \dd\tau_3\,
    \braket{T_C u_1 u_2 u_a u_b u_c u_d}_0.
    \label{eq:first_order_quartic_green_correction}
\end{equation}
The expectation value in Eq.~\eqref{eq:first_order_quartic_green_correction} contains six fields, so it is not forced to vanish. Wick's theorem can pair two of those fields with the external operators and the remaining four fields into two internal contractions. The Hartree-type term is the one in which two fields remain as a correlator at the same contour time and the other two are attached to the external legs. This gives a two-point correction proportional to an equal-time correlator.

A more transparent operator derivation comes from direct mean-field decoupling. Start from the quartic operator product
\begin{equation}
    u_i u_j u_k u_l.
\end{equation}
Use Wick-style pair decoupling,
\begin{align}
    u_i u_j u_k u_l
    \approx
    &\braket{u_i u_j}u_k u_l
    + \braket{u_i u_k}u_j u_l
    + \braket{u_i u_l}u_j u_k \\
    &+ \braket{u_j u_k}u_i u_l
    + \braket{u_j u_l}u_i u_k
    + \braket{u_k u_l}u_i u_j
    - \text{constant terms}.
    \label{eq:quartic_pair_decoupling_again}
\end{align}
Substitute Eq.~\eqref{eq:quartic_pair_decoupling_again} into Eq.~\eqref{eq:quartic_hamiltonian_general_again}. Because $\Phi^{(4)}_{ijkl}$ is fully symmetric, all six pairings lead to the same effective two-point kernel after relabeling dummy indices. Hence the induced quadratic Hamiltonian is
\begin{equation}
    H^{(4)}_{\rm MF}
    =
    \frac{1}{2}\sum_{ij}\Delta K^{(4)}_{ij}u_i u_j + \const,
\end{equation}
with
\begin{equation}
    \Delta K^{(4)}_{ij}
    =
    \frac{1}{2}\sum_{kl}\Phi^{(4)}_{ijkl}\braket{u_k u_l}
    \times \text{(symmetry factor depending on convention)}.
    \label{eq:quartic_deltaK_from_decoupling}
\end{equation}
Equation~\eqref{eq:quartic_deltaK_from_decoupling} is the origin of quartic Hartree mean field.

To rewrite Eq.~\eqref{eq:quartic_deltaK_from_decoupling} in NEGF language, use the equal-time lesser Green's function. By definition,
\begin{equation}
    G^<_{kl}(t,t) = -\frac{\ii}{\hbar}\braket{u_l(t)u_k(t)}.
\end{equation}
Since equal-time displacements commute,
\begin{equation}
    \braket{u_k(t)u_l(t)} = \braket{u_l(t)u_k(t)} = \ii\hbar G^<_{kl}(t,t).
    \label{eq:equal_time_displacement_from_lesser_again}
\end{equation}
Substitute Eq.~\eqref{eq:equal_time_displacement_from_lesser_again} into Eq.~\eqref{eq:quartic_deltaK_from_decoupling}. The quartic mean-field kernel becomes a functional of the equal-time lesser Green's function.

In Wang's tensor convention, the symmetry factor simplifies and the result is
\begin{equation}
    \Sigma_{n,jj'}(\tau,\tau') = 3\ii\hbar\,\delta(\tau,\tau')\sum_{kl}T_{jj'kl}G_{kl}(\tau,\tau).
    \label{eq:wang_quartic_scmf_again}
\end{equation}
The prefactor $3$ comes from the three distinct pairings of a quartic vertex into a two-point external kernel and an internal two-point correlator. This is the basic derivation behind quartic SCMF.

The leading dynamic quartic scattering term appears only at second order in the quartic tensor, i.e. from two quartic vertices. Using Eq.~\eqref{eq:internal_line_count_general} with $m=4$ and $V_4=2$ gives
\begin{equation}
    N_G = \frac{4\times 2 - 2}{2} = 3.
\end{equation}
Therefore the leading dynamic quartic self-energy contains three internal propagators. This is why quartic SCBA is substantially heavier numerically than quartic mean field.

\subsection{What should be called a general \texorpdfstring{$n$}{n}-phonon process in code?}
The derivations above point to a precise code-level abstraction.
\begin{enumerate}[label=\arabic*.]
\item A material input should be represented by the harmonic matrix $K$ and the anharmonic tensors $\Phi^{(3)}, \Phi^{(4)}, \ldots$.
\item An approximation scheme should be defined by which contractions of those tensors with Green's functions are retained.
\item A toy-model benchmark should be a special sparse choice of those tensors, together with a paper-specific unit convention and observable definition.
\end{enumerate}
Thus, in a future materials workflow, ``three-phonon LO'' should mean: start from $\Phi^{(3)}$, build the two-vertex cubic self-energy, evaluate it with the chosen internal propagators, and insert it into Dyson/Keldysh. Likewise, ``four-phonon MF'' should mean: start from $\Phi^{(4)}$, build the quartic Hartree contraction, and solve the resulting effective harmonic problem self-consistently. This tensor-first interpretation is the theory reason the codebase should remain general and approximation-first.

\section{Approximation Schemes: LO, MF, and SCBA}

\subsection{Why multiple approximation schemes are unavoidable}
The exact interacting central-region problem is defined by the coupled equations
\begin{align}
    G^r &= \sqbrak{(G_0^r)^{-1} - \Sigma_n^r[G]}^{-1},
    \label{eq:exact_dyson_closure_repeat} \\
    G^< &= G^r\paren{\Sigma_L^< + \Sigma_R^< + \Sigma_n^<[G]}G^a.
    \label{eq:exact_keldysh_closure_repeat}
\end{align}
Equation~\eqref{eq:exact_dyson_closure_repeat} is exact, but it is not yet useful computationally because the nonlinear self-energy depends on the very Green's function we are trying to solve for. The same statement applies to Eq.~\eqref{eq:exact_keldysh_closure_repeat}. Therefore the central difficulty of anharmonic NEGF is not writing Dyson and Keldysh equations, but closing them.

To make that precise, let us formally expand the nonlinear self-energy in diagram classes. Symbolically,
\begin{equation}
    \Sigma_n[G]
    =
    \Sigma_{\rm cubic}[G] + \Sigma_{\rm quartic}[G] + \cdots.
    \label{eq:self_energy_split_by_interaction_order}
\end{equation}
Each contribution on the right-hand side is itself an infinite sum of diagrams,
\begin{equation}
    \Sigma_{\rm cubic}[G] = \Sigma_{\rm cubic}^{(2v)}[G] + \Sigma_{\rm cubic}^{(4v)}[G] + \cdots,
\end{equation}
\begin{equation}
    \Sigma_{\rm quartic}[G] = \Sigma_{\rm quartic}^{(1v)}[G] + \Sigma_{\rm quartic}^{(2v)}[G] + \cdots,
\end{equation}
where the superscript indicates the number of interaction vertices retained in the diagrammatic expansion. Thus an approximation scheme must answer two questions.
\begin{enumerate}[label=\arabic*.]
\item Which diagram topologies are kept?
\item When those diagrams are evaluated, are the internal lines harmonic propagators $G_0$, fully dressed propagators $G$, or some effective partially renormalized propagators?
\end{enumerate}
Lowest order, mean field, and SCBA are different answers to those two questions.

\subsection{Lowest order (LO)}
\subsubsection{Definition}
Lowest order means: choose the first nonzero self-energy diagrams generated by the chosen interaction and evaluate all internal lines with the harmonic propagators $G_0$.

Let the first nonvanishing nonlinear self-energy be denoted by $\Sigma_n^{[p]}$, where $p$ is the interaction order in the perturbation series. Then LO means
\begin{equation}
    \Sigma_n^{\rm LO} = \Sigma_n^{[p]}[G_0].
    \label{eq:LO_definition_precise}
\end{equation}
For cubic interactions one typically has $p=2$ because the leading dynamic self-energy comes from two cubic vertices. For quartic Hartree mean field, the first nonzero contribution already appears at one quartic vertex.

To derive how LO enters the Green's function, start from the exact Dyson equation in integral form,
\begin{equation}
    G^r = G_0^r + G_0^r\Sigma_n^r G^r.
    \label{eq:dyson_integral_for_LO_derivation}
\end{equation}
Substitute the right-hand side recursively for $G^r$ itself:
\begin{align}
    G^r
    &= G_0^r + G_0^r\Sigma_n^r\paren{G_0^r + G_0^r\Sigma_n^r G^r} \\
    &= G_0^r + G_0^r\Sigma_n^r G_0^r + G_0^r\Sigma_n^r G_0^r\Sigma_n^r G^r.
\end{align}
Continuing this recursion generates the Neumann series
\begin{equation}
    G^r = G_0^r + G_0^r\Sigma_n^r G_0^r + G_0^r\Sigma_n^r G_0^r\Sigma_n^r G_0^r + \cdots.
    \label{eq:dyson_neumann_series_again}
\end{equation}
If $\Sigma_n^r$ is already small in the perturbative sense, the first correction to $G_0^r$ is the term $G_0^r\Sigma_n^r G_0^r$. Thus the LO retarded Green's function is
\begin{equation}
    G^{r,\rm LO}
    \approx
    G_0^r + G_0^r\Sigma_n^{r,\rm LO}[G_0]G_0^r.
    \label{eq:LO_retarded_from_neumann}
\end{equation}
We can derive the corresponding lesser formula in exactly the same spirit. Start from the exact Keldysh equation,
\begin{equation}
    G^< = G^r\Sigma_{\rm tot}^< G^a,
    \qquad
    \Sigma_{\rm tot}^< = \Sigma_L^< + \Sigma_R^< + \Sigma_n^<.
    \label{eq:exact_keldysh_for_LO_derivation}
\end{equation}
Write
\begin{equation}
    G^r = G_0^r + \delta G^r,
    \qquad
    G^a = G_0^a + \delta G^a,
    \qquad
    \Sigma_{\rm tot}^< = \Sigma_0^< + \Sigma_n^<,
\end{equation}
with
\begin{equation}
    \Sigma_0^< = \Sigma_L^< + \Sigma_R^<.
\end{equation}
Retaining only terms linear in the nonlinear self-energy gives
\begin{align}
    G^<
    &\approx
    G_0^r\Sigma_0^<G_0^a
    + \delta G^r\Sigma_0^<G_0^a
    + G_0^r\Sigma_0^<\delta G^a
    + G_0^r\Sigma_n^<G_0^a.
\end{align}
Now use
\begin{equation}
    G_0^< = G_0^r\Sigma_0^<G_0^a,
\end{equation}
plus
\begin{equation}
    \delta G^r = G_0^r\Sigma_n^{r,\rm LO}G_0^r,
    \qquad
    \delta G^a = G_0^a\Sigma_n^{a,\rm LO}G_0^a.
\end{equation}
Then
\begin{align}
    G^{<,\rm LO}
    &\approx
    G_0^<
    + G_0^r\Sigma_n^{r,\rm LO}G_0^r\Sigma_0^<G_0^a
    + G_0^r\Sigma_0^<G_0^a\Sigma_n^{a,\rm LO}G_0^a
    + G_0^r\Sigma_n^{<,\rm LO}G_0^a \\
    &=
    G_0^<
    + G_0^r\Sigma_n^{r,\rm LO}G_0^<
    + G_0^<\Sigma_n^{a,\rm LO}G_0^a
    + G_0^r\Sigma_n^{<,\rm LO}G_0^a.
    \label{eq:LO_lesser_full_derivation}
\end{align}
Equation~\eqref{eq:LO_lesser_full_derivation} is the explicit linearized Keldysh expansion behind a lowest-order inelastic conductance formula.

\subsubsection{Assumptions}
The assumptions of LO are now transparent from the derivation.
\begin{enumerate}[label=\arabic*.]
\item The retained self-energy $\Sigma_n^{\rm LO}$ is small enough that terms quadratic or higher in $\Sigma_n$ can be dropped in Eq.~\eqref{eq:dyson_neumann_series_again}.
\item The internal lines inside the self-energy diagrams are still well approximated by harmonic propagators.
\item Observables can be trusted even though the feedback of scattering on the internal lines is neglected.
\item The omitted higher-topology diagrams are genuinely subleading in the regime of interest.
\end{enumerate}
Because LO is a truncation before self-consistency, it is generally not a conserving approximation in the strict Baym--Kadanoff sense. It is nevertheless the clearest perturbative benchmark.

\subsubsection{Strengths and limitations}
The strengths of LO follow directly from Eq.~\eqref{eq:LO_retarded_from_neumann} and Eq.~\eqref{eq:LO_lesser_full_derivation}.
\begin{itemize}
\item Every retained correction can be traced to a definite diagram and a definite interaction order.
\item The approximation is analytically transparent and often admits closed-form toy-model benchmarks.
\item Weak-coupling lifetimes derived from LO connect naturally to Fermi's golden rule.
\end{itemize}
The limitations also follow directly from the same equations.
\begin{itemize}
\item Once $\Sigma_n$ is not small, the neglected quadratic and higher terms in Eq.~\eqref{eq:dyson_neumann_series_again} matter.
\item The harmonic internal lines may become a poor representation of the actual broadened spectrum.
\item Conservation relations can be degraded because the self-energy and Green's functions are not solved self-consistently.
\end{itemize}

\subsubsection{Examples from the local paper set}
The local paper set gives three canonical realizations of LO.
\begin{enumerate}[label=\arabic*.]
\item Wang 2007 Fig.~5: a cubic onsite toy model in which the first nontrivial cubic graph is evaluated explicitly.
\item Xu \emph{et al.} 2008: the retarded cubic self-energy is evaluated on harmonic modes to extract linewidths and lifetimes.
\item Wang 2006: the nonlinear correction enters perturbatively in the center region of a mesoscopic junction calculation.
\end{enumerate}
In all three cases the approximation logic is the same: identify the first nonzero diagram class and evaluate it with the harmonic propagators.

\subsection{Mean field (MF)}
\subsubsection{Definition}
Mean field replaces higher-order operator products by lower-order expectation values, thereby reducing the interacting Hamiltonian to an effective quadratic Hamiltonian. The formal pattern is
\begin{equation}
    u_{i_1}\cdots u_{i_m}
    \longrightarrow
    \sum \braket{u\cdots u}\times \text{lower-order operator products}.
    \label{eq:MF_operator_pattern}
\end{equation}
The resulting effective Hamiltonian is harmonic but depends on correlators that must be determined self-consistently.

The cleanest derivation is the quartic case. Start from
\begin{equation}
    H^{(4)} = \frac{1}{4!}\sum_{ijkl}\Phi^{(4)}_{ijkl}u_i u_j u_k u_l.
    \label{eq:quartic_start_for_MF}
\end{equation}
Use the pair decoupling
\begin{align}
    u_i u_j u_k u_l
    \approx
    &\braket{u_i u_j}u_k u_l
    + \braket{u_i u_k}u_j u_l
    + \braket{u_i u_l}u_j u_k \\
    &+ \braket{u_j u_k}u_i u_l
    + \braket{u_j u_l}u_i u_k
    + \braket{u_k u_l}u_i u_j
    - \text{constant terms}.
    \label{eq:quartic_MF_decoupling_repeat}
\end{align}
Substitute Eq.~\eqref{eq:quartic_MF_decoupling_repeat} into Eq.~\eqref{eq:quartic_start_for_MF}. Because the quartic tensor is symmetric, the six contractions produce one common quadratic structure after relabeling dummy indices. Hence
\begin{equation}
    H^{(4)}_{\rm MF}
    =
    \frac{1}{2}\sum_{ij}\Delta K^{(4)}_{ij}u_i u_j + \const.
    \label{eq:quartic_MF_effective_hamiltonian}
\end{equation}
The effective renormalization matrix can be written as
\begin{equation}
    \Delta K^{(4)}_{ij}
    =
    C^{(4)}\sum_{kl}\Phi^{(4)}_{ijkl}\braket{u_k u_l},
    \label{eq:quartic_MF_deltaK_general}
\end{equation}
where $C^{(4)}$ is the convention-dependent combinatorial factor.

Therefore the full effective harmonic matrix is
\begin{equation}
    K_{\rm eff} = K^C + \Delta K^{(4)}.
    \label{eq:effective_harmonic_matrix_MF}
\end{equation}
Once Eq.~\eqref{eq:effective_harmonic_matrix_MF} is known, the effective retarded Green's function is simply
\begin{equation}
    G^r_{\rm MF}(\omega)
    =
    \sqbrak{(\omega+\ii\eta)^2I - K_{\rm eff} - \Sigma_L^r - \Sigma_R^r}^{-1}.
    \label{eq:MF_retarded_green_from_keff}
\end{equation}
The key point is that $K_{\rm eff}$ depends on the equal-time correlator, which depends on the lesser Green's function of the effective problem. Thus mean field is a self-consistent effective harmonic theory.

To connect Eq.~\eqref{eq:quartic_MF_deltaK_general} to NEGF notation, use the identity
\begin{equation}
    G^<_{kl}(t,t) = -\frac{\ii}{\hbar}\braket{u_l(t)u_k(t)}.
\end{equation}
Since equal-time displacement operators commute, we may write
\begin{equation}
    \braket{u_k u_l} = \ii\hbar G^<_{kl}(t,t).
    \label{eq:equal_time_correlation_to_lesser_MF}
\end{equation}
Substituting Eq.~\eqref{eq:equal_time_correlation_to_lesser_MF} into Eq.~\eqref{eq:quartic_MF_deltaK_general} gives a self-energy functional of the form
\begin{equation}
    \Sigma_n^{\rm MF}[G]
    =
    \Sigma_n^{\rm static}[G(t,t)].
    \label{eq:MF_self_energy_functional_form}
\end{equation}
In Wang's quartic tensor convention, Eq.~\eqref{eq:MF_self_energy_functional_form} becomes
\begin{equation}
    \Sigma_{n,jj'}(\tau,\tau') = 3\ii\hbar\,\delta(\tau,\tau')\sum_{kl}T_{jj'kl}G_{kl}(\tau,\tau).
    \label{eq:wang_quartic_MF_closure_repeat}
\end{equation}
Equation~\eqref{eq:wang_quartic_MF_closure_repeat} is therefore not an ad hoc closure; it is the quartic Hartree contraction written in the Wang tensor convention.

\subsubsection{Assumptions}
The assumptions of MF follow directly from Eq.~\eqref{eq:quartic_MF_effective_hamiltonian} and Eq.~\eqref{eq:MF_self_energy_functional_form}.
\begin{enumerate}[label=\arabic*.]
\item Static renormalization is assumed to dominate over dynamic broadening.
\item Higher-order operator products are assumed to be well approximated by pair correlators.
\item The main effect of the interaction is to renormalize the effective force constants or frequencies.
\item The omitted dynamic scattering channels are assumed to be less important than the retained static contractions.
\end{enumerate}
Mean field is therefore most natural when the interaction mainly shifts the spectrum rather than producing strong inelastic linewidths.

\subsubsection{Three-phonon MF}
Cubic interactions can also be decoupled in a mean-field-like way, but the resulting structure is less natural. Start from
\begin{equation}
    H^{(3)} = \frac{1}{3!}\sum_{ijk}\Phi^{(3)}_{ijk}u_i u_j u_k.
    \label{eq:cubic_start_for_MF}
\end{equation}
A naive decoupling gives
\begin{equation}
    u_i u_j u_k
    \approx
    \braket{u_i u_j}u_k + \braket{u_i u_k}u_j + \braket{u_j u_k}u_i
    + \braket{u_i}u_j u_k + \braket{u_j}u_i u_k + \braket{u_k}u_i u_j.
    \label{eq:cubic_MF_decoupling_explicit}
\end{equation}
If the reference state is expanded around a centered equilibrium point, then
\begin{equation}
    \braket{u_i}=0.
\end{equation}
Thus the cubic mean-field correction reduces mainly to a tadpole-like effective linear term proportional to $\braket{u_i u_j}u_k$. Such a term often indicates a shift of the reference equilibrium configuration rather than a faithful description of transport scattering. This is why cubic interactions are usually benchmarked first through LO or SCBA rather than through MF.

\subsubsection{Four-phonon MF}
Quartic interactions are the natural domain of MF because Eq.~\eqref{eq:quartic_MF_effective_hamiltonian} is already quadratic. The self-consistent algorithm follows directly from the equations.
\begin{enumerate}[label=\arabic*.]
\item Start from an initial guess for the equal-time correlator $\braket{u_k u_l}$.
\item Construct $\Delta K^{(4)}$ from Eq.~\eqref{eq:quartic_MF_deltaK_general}.
\item Form $K_{\rm eff}$ from Eq.~\eqref{eq:effective_harmonic_matrix_MF}.
\item Solve for $G^r_{\rm MF}$ from Eq.~\eqref{eq:MF_retarded_green_from_keff}.
\item Compute $G^<_{\rm MF}$ from the Keldysh equation of the effective harmonic problem.
\item Update the equal-time correlator using Eq.~\eqref{eq:equal_time_correlation_to_lesser_MF}.
\item Repeat until $K_{\rm eff}$ or the correlator converges.
\end{enumerate}
This is precisely the logic of self-consistent mean field (SCMF).

\subsubsection{Strengths and limitations}
The strengths of MF are now explicit.
\begin{itemize}
\item It directly captures quartic Hartree renormalization.
\item It reduces the nonlinear problem to a self-consistent harmonic problem.
\item It is often numerically more stable than a fully dynamic self-consistency.
\end{itemize}
Its limitations are equally explicit.
\begin{itemize}
\item It neglects genuine dynamic scattering and linewidth broadening.
\item It may misrepresent transport when the dominant physics is inelastic relaxation rather than static renormalization.
\item For cubic interactions it is usually not the most natural physical approximation.
\end{itemize}

\subsection{Self-consistent Born approximation (SCBA)}
\subsubsection{Definition}
SCBA keeps the Born-level diagram topology but replaces the harmonic internal propagators by dressed propagators determined self-consistently. The difference from LO is therefore not in topology, but in the choice of internal lines.

Suppose the Born self-energy functional is
\begin{equation}
    \Sigma_n^{\rm Born}[X].
\end{equation}
Then LO uses
\begin{equation}
    \Sigma_n^{\rm LO} = \Sigma_n^{\rm Born}[G_0],
    \label{eq:LO_as_Born_on_G0}
\end{equation}
whereas SCBA uses
\begin{equation}
    \Sigma_n^{\rm SCBA} = \Sigma_n^{\rm Born}[G].
    \label{eq:SCBA_as_Born_on_G}
\end{equation}
For example, in a cubic problem the Born structure is quadratic in the internal propagators. Thus,
\begin{equation}
    \Sigma_{\rm cubic}^{\rm Born}[X] \propto \Phi^{(3)} X X \Phi^{(3)}.
\end{equation}
Replacing $X=G_0$ gives LO, while replacing $X=G$ gives SCBA.

Insert Eq.~\eqref{eq:SCBA_as_Born_on_G} into Dyson's equation:
\begin{equation}
    G^r = \sqbrak{(G_0^r)^{-1} - \Sigma_n^{r,\rm Born}[G]}^{-1}.
    \label{eq:SCBA_retarded_fixed_point}
\end{equation}
Likewise, for the lesser sector,
\begin{equation}
    G^< = G^r\paren{\Sigma_L^< + \Sigma_R^< + \Sigma_n^{<,\rm Born}[G]}G^a.
    \label{eq:SCBA_lesser_fixed_point}
\end{equation}
Equation~\eqref{eq:SCBA_retarded_fixed_point} and Eq.~\eqref{eq:SCBA_lesser_fixed_point} define a nonlinear fixed-point problem.

We can derive the iterative structure explicitly. Start from an initial guess
\begin{equation}
    G^{r,(0)} = G_0^r,
    \qquad
    G^{<,(0)} = G_0^<.
    \label{eq:SCBA_initial_guess}
\end{equation}
At iteration $n$, evaluate
\begin{equation}
    \Sigma_n^{r,(n+1)} = \Sigma_n^{r,\rm Born}[G^{(n)}],
    \qquad
    \Sigma_n^{<,(n+1)} = \Sigma_n^{<,\rm Born}[G^{(n)}].
    \label{eq:SCBA_self_energy_update}
\end{equation}
Then update the Green's functions using
\begin{equation}
    G^{r,(n+1)} = \sqbrak{(G_0^r)^{-1} - \Sigma_n^{r,(n+1)}}^{-1},
    \label{eq:SCBA_retarded_update}
\end{equation}
\begin{equation}
    G^{<,(n+1)} = G^{r,(n+1)}\paren{\Sigma_L^< + \Sigma_R^< + \Sigma_n^{<,(n+1)}}G^{a,(n+1)}.
    \label{eq:SCBA_lesser_update}
\end{equation}
The iteration continues until a convergence criterion is met, for example
\begin{equation}
    \norm{G^{r,(n+1)} - G^{r,(n)}} < \varepsilon,
    \qquad
    \norm{G^{<,(n+1)} - G^{<,(n)}} < \varepsilon.
    \label{eq:SCBA_convergence_criterion}
\end{equation}
This is the explicit fixed-point realization of SCBA.

\subsubsection{Assumptions}
The assumptions of SCBA follow directly from Eq.~\eqref{eq:SCBA_as_Born_on_G}.
\begin{enumerate}[label=\arabic*.]
\item The dominant missing physics in LO is the feedback of scattering on the internal propagators.
\item The retained Born topology is assumed to be more important than omitted higher-topology vertex corrections.
\item The self-consistent fixed point is assumed to exist and to represent the physically relevant branch.
\end{enumerate}
Thus SCBA is not ``exact,'' but it is a systematic resummation of the chosen Born diagram class.

\subsubsection{Why SCBA is attractive}
The attraction of SCBA is easiest to see by comparing it with the LO Dyson series. In LO,
\begin{equation}
    G^{r,\rm LO} \approx G_0^r + G_0^r\Sigma_n^{r,\rm LO}[G_0]G_0^r,
\end{equation}
so the internal lines never feel the broadening produced by the self-energy. In SCBA, by contrast, the updated Green's function is immediately recycled into the next self-energy evaluation through Eq.~\eqref{eq:SCBA_self_energy_update}. Therefore broadening, spectral shifts, and nonequilibrium populations feed back onto the scattering kernel itself. This is often essential once the interaction is no longer a tiny perturbation.

\subsubsection{Why SCBA is difficult}
The difficulty of SCBA is also explicit in the equations. Every iteration requires:
\begin{enumerate}[label=\arabic*.]
\item evaluating the retarded nonlinear self-energy,
\item evaluating the lesser nonlinear self-energy,
\item solving Dyson's equation,
\item solving Keldysh's equation,
\item checking convergence and possibly applying mixing or stabilization.
\end{enumerate}
In realistic phonon problems the self-energy evaluation itself contains frequency convolutions. This is why SCBA is substantially more expensive than LO or MF, and why many codes first implement simpler SCBA-like closures before attempting the full contour-resolved version.

\subsection{Comparison table}
\begin{center}
\begin{longtable}{p{0.14\linewidth}p{0.22\linewidth}p{0.24\linewidth}p{0.28\linewidth}}
\toprule
Scheme & Internal lines & Typical physics captured & Main risk \\
\midrule
LO & $G_0$ & first correction, perturbative linewidth, benchmark clarity & misses feedback and may fail at moderate anharmonicity \\
MF & effective static $G$ from renormalized $K$ & frequency renormalization, static quartic effects & misses dynamic broadening and true inelastic channels \\
SCBA & dressed $G$ & feedback of scattering and renormalization, conserving tendencies & numerically expensive and still diagram-class dependent \\
\bottomrule
\end{longtable}
\end{center}
The table is a short summary, but the equation-level distinctions are now clear:
\begin{enumerate}[label=\arabic*.]
\item LO truncates both topology and internal-line feedback.
\item MF keeps a static contraction class and solves the resulting effective harmonic problem self-consistently.
\item SCBA keeps a dynamic Born topology and restores feedback by dressing its internal lines self-consistently.
\end{enumerate}

\subsection{Practical guidance for our code design}
The derivations above translate directly into software design rules.
\begin{enumerate}[label=\arabic*.]
\item The reusable inputs should be the harmonic matrix $K$, the anharmonic tensors $\Phi^{(3)}, \Phi^{(4)}, \ldots$, and the lead self-energies.
\item LO, MF, and SCBA should remain separate modules because each one corresponds to a different closure of the exact Dyson--Keldysh system.
\item Toy models should remain thin wrappers that provide special sparse tensors and paper-specific conventions.
\item The material-ready path should inherit the exact tensor-level derivations, not the toy-model shortcuts used in figure benchmarks.
\end{enumerate}
This separation is the theory basis for keeping the code approximation-first and tensor-first.

\section{Recovery of the Main Formulas in the Local Literature Set}

\subsection{Guide to this section}
The purpose of this section is to recover, equation by equation, the main formulas used in the papers stored in the local literature set. The word ``recover'' is important: rather than merely quoting a formula, we want to show exactly how it follows from the general phonon-NEGF derivations in the earlier sections.

This serves two functions.
\begin{enumerate}[label=\arabic*.]
\item It checks that the general theory note is genuinely consistent with the literature benchmarks.
\item It identifies which steps in a paper are general theory and which steps are special toy-model assumptions or notation choices.
\end{enumerate}
The section is therefore organized around derivation chains rather than around citation summaries.

\subsection{Paper set overview}
\begin{longtable}{p{0.19\linewidth}p{0.31\linewidth}p{0.41\linewidth}}
\toprule
Local file & Title & Main formulas we want to recover \\
\midrule
\texttt{0605028v1.pdf} & Wang--Wang--Zeng 2006, NEGF approach to mesoscopic thermal transport & current formula, perturbative nonlinear self-energy usage \\
\texttt{0701164v1.pdf} & Wang \emph{et al.} 2007, thermal transport in junctions & Dyson/Keldysh structure, effective transmission, cubic LO toy model, SCMF logic \\
\texttt{0802.2761v1.pdf} & Wang--Wang--Lu 2008 review & Landauer, Caroli, surface Green's functions, contour formalism summary \\
\texttt{0807.3819v1.pdf} & Xu \emph{et al.} 2008, phonon-phonon interaction and ballistic-diffusive transport & momentum-space retarded self-energy, lifetime, mean free path, phenomenological transmission \\
\texttt{1303.7317v1.pdf} & Wang \emph{et al.} 2013 review & general NEGF review, nonlinear comments, quartic SCMF closure \\
\texttt{transp.pdf} & journal version of the 2008 review & same core formulas as \texttt{0802.2761v1.pdf} \\
\bottomrule
\end{longtable}

\subsection{Recovery of Wang 2006 (\texttt{0605028v1.pdf})}
The 2006 paper already contains the basic lead--center--lead contour-NEGF structure. We recover its main formulas in three steps: eliminate the leads, write the current, then add the nonlinear self-energy.

Start from the harmonic block Hamiltonian and the block matrix of Eq.~\eqref{eq:blockK}. Eliminating the leads gave us
\begin{equation}
    G_0^r(\omega)
    =
    \sqbrak{(\omega+\ii\eta)^2I - K^C - \Sigma_L^r(\omega) - \Sigma_R^r(\omega)}^{-1}.
    \label{eq:recover_2006_G0r}
\end{equation}
This is precisely the central harmonic Green's function used in Wang 2006, though not necessarily in identical notation.

Next recover the current formula. From the foundations derivation we obtained
\begin{equation}
    I_L = -\int \frac{\dd\omega}{2\pi}\,\hbar\omega\,\Tr\sqbrak{G^r\Sigma_L^< + G^<\Sigma_L^a}.
    \label{eq:recover_2006_current_repeat}
\end{equation}
Equation~\eqref{eq:recover_2006_current_repeat} is already the central transport formula of the 2006 paper. The only remaining step is to understand how nonlinearity is inserted.

Let the nonlinear center self-energy be $\Sigma_n$. Then the exact interacting equations are
\begin{align}
    G^r &= \sqbrak{(G_0^r)^{-1} - \Sigma_n^r}^{-1}, \\
    G^< &= G^r\paren{\Sigma_L^< + \Sigma_R^< + \Sigma_n^<}G^a.
    \label{eq:recover_2006_interacting_pair}
\end{align}
Thus Wang 2006 may be viewed as the first paper in the local set that already contains the full modern logic:
\begin{equation}
    \text{harmonic lead elimination}
    \Rightarrow
    G_0^r
    \Rightarrow
    \Sigma_n
    \Rightarrow
    G^r,G^<
    \Rightarrow
    I_L.
\end{equation}
The later papers then refine specific parts of that chain rather than replacing it.

\subsection{Recovery of Wang 2007 (\texttt{0701164v1.pdf})}
This paper is the key literature benchmark for the current project. We recover its main formulas one by one.

\subsubsection{Hamiltonian structure}
The paper writes the nonlinear Hamiltonian as
\begin{equation}
    H_n = \frac{1}{3}\sum_{ijk}T_{ijk}u_i u_j u_k + \frac{1}{4}\sum_{ijkl}T_{ijkl}u_i u_j u_k u_l.
    \label{eq:wang2007_Hn_repeated_again}
\end{equation}
To connect this with the symmetric IFC convention, start from
\begin{equation}
    H_n = \frac{1}{3!}\sum_{ijk}\Phi^{(3)}_{ijk}u_i u_j u_k + \frac{1}{4!}\sum_{ijkl}\Phi^{(4)}_{ijkl}u_i u_j u_k u_l.
    \label{eq:symmetric_Hn_for_wang2007_recovery}
\end{equation}
Compare cubic terms:
\begin{equation}
    \frac{1}{3}T_{ijk}u_i u_j u_k = \frac{1}{3!}\Phi^{(3)}_{ijk}u_i u_j u_k.
\end{equation}
Multiply by $6$:
\begin{equation}
    2T_{ijk}u_i u_j u_k = \Phi^{(3)}_{ijk}u_i u_j u_k.
\end{equation}
Hence
\begin{equation}
    \Phi^{(3)}_{ijk} = 2T_{ijk}.
    \label{eq:wang2007_phi3_conversion}
\end{equation}
Likewise for the quartic term,
\begin{equation}
    \frac{1}{4}T_{ijkl}u_i u_j u_k u_l = \frac{1}{4!}\Phi^{(4)}_{ijkl}u_i u_j u_k u_l,
\end{equation}
so after multiplying by $24$,
\begin{equation}
    6T_{ijkl}u_i u_j u_k u_l = \Phi^{(4)}_{ijkl}u_i u_j u_k u_l,
\end{equation}
and therefore
\begin{equation}
    \Phi^{(4)}_{ijkl} = 6T_{ijkl}.
    \label{eq:wang2007_phi4_conversion}
\end{equation}
So the Wang 2007 Hamiltonian is exactly the general IFC Hamiltonian with a different vertex normalization.

\subsubsection{Dyson and Keldysh equations}
The paper's contour equations reduce to the same central algebraic equations already derived generally. Starting from the exact center Dyson equation,
\begin{equation}
    G^r = G_0^r + G_0^r\Sigma_n^r G^r,
\end{equation}
we derived
\begin{equation}
    G^r = \sqbrak{(G_0^r)^{-1} - \Sigma_n^r}^{-1}.
    \label{eq:wang2007_recovered_dyson}
\end{equation}
Similarly, the lesser equation is
\begin{equation}
    G^< = G^r\paren{\Sigma^< + \Sigma_n^<}G^a,
    \qquad
    \Sigma^< = \Sigma_L^< + \Sigma_R^<.
    \label{eq:wang2007_recovered_keldysh}
\end{equation}
Thus the formal backbone of Wang 2007 is not a different formalism from ours; it is the same Dyson--Keldysh structure, specialized to explicit benchmark calculations.

\subsubsection{Current formula}
The paper's current expression is recovered directly from the general current formula,
\begin{equation}
    I_L = -\int \frac{\dd\omega}{2\pi}\,\hbar\omega\,\Tr\sqbrak{G^r\Sigma_L^< + G^<\Sigma_L^a}.
    \label{eq:wang2007_recovered_current}
\end{equation}
This should be read carefully: in the interacting theory, Eq.~\eqref{eq:wang2007_recovered_current} is exact if the exact $G^r$ and $G^<$ are used. All approximation errors are therefore pushed into the approximation chosen for $\Sigma_n$.

\subsubsection{Cubic onsite model formulas}
The paper's Fig.~5 cubic benchmark is a nearest-neighbor harmonic chain with onsite harmonic confinement and a cubic onsite interaction. The equation of motion is
\begin{equation}
    \ddot u_j = K u_{j-1} + (-2K-K_0)u_j + K u_{j+1} - t_j u_j^2.
    \label{eq:wang2007_onsite_eom_again}
\end{equation}
To recover the harmonic Green's function, first remove the nonlinear term. The harmonic recurrence relation is then
\begin{equation}
    K u_{j-1} + (-2K-K_0)u_j + K u_{j+1} = -\ddot u_j.
\end{equation}
Fourier transform in time using $\ddot u_j \to -(\omega+\ii\eta)^2u_j$ for the retarded solution. Then
\begin{equation}
    K u_{j-1} + \Omega u_j + K u_{j+1} = 0,
    \qquad
    \Omega = (\omega+\ii\eta)^2 - 2K - K_0.
    \label{eq:wang2007_harmonic_recurrence}
\end{equation}
Seek a solution of the form
\begin{equation}
    u_j \propto \lambda^j.
    \label{eq:lambda_ansatz}
\end{equation}
Substitute Eq.~\eqref{eq:lambda_ansatz} into Eq.~\eqref{eq:wang2007_harmonic_recurrence}:
\begin{equation}
    K\lambda^{j-1} + \Omega \lambda^j + K\lambda^{j+1} = 0.
\end{equation}
Divide by $\lambda^j$:
\begin{equation}
    K\lambda^{-1} + \Omega + K\lambda = 0.
    \label{eq:lambda_characteristic_equation_again}
\end{equation}
This is exactly the quadratic equation quoted in the paper. Solving Eq.~\eqref{eq:lambda_characteristic_equation_again} gives the two roots $\lambda_1$ and $\lambda_2$. The retarded boundary condition selects the root with the appropriate decay property, and the Green's function becomes
\begin{equation}
    G_{0,jl}^r(\omega) = \frac{\lambda_1^{\abs{j-l}}}{(\lambda_1-\lambda_2)K}.
    \label{eq:wang2007_harmonic_green_recovered}
\end{equation}
Now recover Eq.~(77), the first cubic graph. The onsite interaction is
\begin{equation}
    H_n = \frac{1}{3}\sum_j t_j u_j^3.
    \label{eq:onsite_cubic_hamiltonian_for_eq77}
\end{equation}
Expand the contour Green's function to second order in $H_n$:
\begin{equation}
    G = G_0 + \frac{1}{2!}\paren{-\frac{\ii}{\hbar}}^2\int_C \dd\tau_1\dd\tau_2\,\braket{T_C u u H_n(\tau_1)H_n(\tau_2)}_0 + \cdots.
    \label{eq:eq77_expansion_start}
\end{equation}
Insert Eq.~\eqref{eq:onsite_cubic_hamiltonian_for_eq77} into Eq.~\eqref{eq:eq77_expansion_start}:
\begin{align}
    G
    =
    G_0
    + \frac{1}{2!}\paren{-\frac{\ii}{\hbar}}^2
      \frac{1}{3^2}
      \sum_{j_1 j_2}t_{j_1}t_{j_2}
      \int_C \dd\tau_1\dd\tau_2\,
      \braket{T_C u u u_{j_1}^3(\tau_1)u_{j_2}^3(\tau_2)}_0
    + \cdots.
    \label{eq:eq77_expansion_after_insertion}
\end{align}
Now apply Wick's theorem. The one-particle-irreducible two-point term is the one where one external line attaches to the first cubic vertex, one external line attaches to the second cubic vertex, and the remaining four internal legs are paired between the two vertices. Because the interaction is local in site index, the tensor structure collapses. After summing the equivalent contractions, the resulting self-energy is
\begin{equation}
    \Sigma_{n,jl}^{\sigma\sigma'}(t) = 2\ii t^2\sqbrak{G_{0,jl}^{\sigma\sigma'}(t)}^2.
    \label{eq:wang2007_eq77_recovered_precise}
\end{equation}
Equation~\eqref{eq:wang2007_eq77_recovered_precise} is exactly the paper's Eq.~(77). The important point is that it is not a special phenomenological ansatz; it is the explicit two-vertex cubic contraction specialized to an onsite scalar interaction.

The paper's second graph is a local tadpole-like contribution. Formally, it comes from the contraction class in which one cubic vertex contributes an equal-time contraction that reduces the remaining structure to a static local correction. This is why, in the code benchmark, the second graph behaves much more like a static shift than like a dynamic broadening term.

Now recover the effective transmission formula. Start from the exact current formula and the corresponding right-lead current expression. In a steady state one has $I_L = -I_R$, so one may symmetrize the current. Insert the interacting $G^r$ and $G^<$ into the exact current and isolate the factor $f_L-f_R$. The remaining matrix object is then defined as
\begin{equation}
    \widetilde{\mathcal T}(\omega)
    =
    \frac{1}{2(f_L-f_R)}
    \Tr\cbrak{(G^r-G^a)(\Sigma_R^< - \Sigma_L^<) + \ii G^<(\Gamma_R-\Gamma_L)}.
    \label{eq:wang2007_effective_transmission_recovered_full}
\end{equation}
This formula is therefore not an independent transport law; it is a convenient rewriting of the exact interacting current into a transmission-like representation.

The paper then takes the linear-response limit explicitly. Let
\begin{equation}
    T_L = T + \frac{\Delta T}{2},
    \qquad
    T_R = T - \frac{\Delta T}{2}.
    \label{eq:wang2007_temperature_split}
\end{equation}
Expand the interacting Green's functions around the equilibrium point $T_L=T_R=T$:
\begin{equation}
    G^r = G_{\rm eq}^r + \frac{\delta G^r}{\delta T}\Delta T + \order{\Delta T^2},
    \qquad
    G^< = G_{\rm eq}^< + \frac{\delta G^<}{\delta T}\Delta T + \order{\Delta T^2}.
    \label{eq:wang2007_eq83_recovered}
\end{equation}
Substitute Eq.~\eqref{eq:wang2007_eq83_recovered} into Eq.~\eqref{eq:wang2007_effective_transmission_recovered_full}. The terms built only from equilibrium Green's functions cancel identically, because the current must vanish when $T_L=T_R$. After dividing by $f_L-f_R$ and taking $\Delta T\to 0$, the remaining first-order object is exactly the paper's effective transmission in linear response.

For the specific one-dimensional chain of Fig.~5, only the first site and the last site couple directly to the left and right baths. Therefore $\Gamma_L$ and $\Gamma_R$ have support only on the boundary sites, and the trace in Eq.~\eqref{eq:wang2007_effective_transmission_recovered_full} collapses to a boundary expression. Using the explicit surface self-energies of the uniform chain, the paper rewrites the result as
\begin{equation}
    \widetilde{\mathcal T}(\omega)
    = iK\,\Im \lambda_1\paren{F_{11}^L - F_{NN}^R},
    \label{eq:wang2007_boundary_F_reduction}
\end{equation}
where the $F$-functions collect the nonequilibrium temperature variations of the Green's functions. The left-boundary function is
\begin{equation}
    F^L
    =
    \frac{1}{2}\paren{G^r-G^a}
    +
    \paren{
        f\frac{\delta(G^r-G^a)}{\delta T}
        -
        \frac{\delta G^<}{\delta T}
    }
    \paren{\frac{\partial f}{\partial T}}^{-1},
    \label{eq:wang2007_eq85_recovered}
\end{equation}
and $F^R$ is the same except that the first term changes sign. Equation~\eqref{eq:wang2007_eq85_recovered} is the paper's Eq.~(85). Its meaning is straightforward: the first term is the equilibrium spectral contribution, while the second term is the explicit linear-response correction due to the temperature dependence of the Green's functions.

The paper then differentiates the Dyson and Keldysh equations with respect to temperature. We recover those formulas directly. Start from
\begin{equation}
    (G^r)^{-1} = (G_0^r)^{-1} - \Sigma_n^r.
\end{equation}
Differentiate with respect to temperature:
\begin{equation}
    \delta (G^r)^{-1} = -\delta\Sigma_n^r.
\end{equation}
Use the matrix identity
\begin{equation}
    \delta G^r = -G^r\,\delta (G^r)^{-1}\,G^r,
\end{equation}
which follows from differentiating $G^r(G^r)^{-1}=I$. Then
\begin{equation}
    \delta G^r = G^r\,\delta\Sigma_n^r\,G^r.
    \label{eq:wang2007_eq86_recovered}
\end{equation}
This is precisely the paper's Eq.~(86). So Eq.~(86) is simply the matrix derivative of Dyson's equation.

The lesser variation is more involved because $G^<$ depends on $G^r$, $G^a$, the lead lesser self-energies, and the nonlinear lesser self-energy. A convenient exact rearrangement is
\begin{equation}
    G^<
    =
    \bar G_0^<
    +
    G^r\Sigma_n^r\bar G_0^<
    +
    G^r\Sigma_n^<G^a,
    \qquad
    \bar G_0^< = G_0^<\paren{I+\Sigma_n^a G^a}.
    \label{eq:wang2007_rearranged_lesser_before_variation}
\end{equation}
Now differentiate Eq.~\eqref{eq:wang2007_rearranged_lesser_before_variation} term by term with respect to temperature. Use the product rule on every factor. Since the paper evaluates the limit at $T_L=T_R=T$, the retarded harmonic Green's function is temperature independent in this toy model and
\begin{equation}
    \delta G_0^r = 0.
    \label{eq:wang2007_deltaG0r_zero}
\end{equation}
After collecting terms, one arrives at the paper's Eq.~(87), which expresses $\delta G^</\delta T$ in terms of $\delta G_0^</\delta T$, $\delta\Sigma_n^r/\delta T$, and $\delta\Sigma_n^</\delta T$. The key structural point is that Eq.~(87) is not a new approximation; it is the exact temperature variation of the interacting lesser Green's function written in a form suitable for the cubic onsite benchmark.

The next quantity needed is the variation of the harmonic lesser Green's function. For the uniform chain, the left- and right-moving parts of the harmonic $G_0^<$ can be written explicitly in terms of the root $\lambda_1$. Differentiating the Bose factors with respect to temperature and then setting $T_L=T_R=T$ yields
\begin{equation}
    \frac{\delta G_{0,jl}^<}{\delta T}
    =
    \frac{\lambda_1^{\,j-l} - \lambda_1^{\,l-j}}{2(\lambda_1-\lambda_2)K}
    \Theta(\omega)\,
    \frac{\partial f}{\partial T}.
    \label{eq:wang2007_eq88_recovered}
\end{equation}
Equation~\eqref{eq:wang2007_eq88_recovered} is the paper's Eq.~(88). Its role is to reduce all temperature dependence to the known derivative of the Bose function, multiplied by a purely harmonic chain kernel.

Finally, recover Eq.~(89). The first cubic graph in frequency space is a convolution of two harmonic Green's functions. Differentiating that convolution with respect to temperature gives two identical terms, because the derivative can act on either internal line. The equality of those two contributions is the origin of the factor of $2$ that upgrades the prefactor from $2it^2$ in the original graph to $4it^2$ in the temperature derivative. The result is
\begin{equation}
    \frac{\delta \Sigma_{n,jl}^{r,<}[\omega]}{\delta T}
    =
    \frac{4it^2}{2\pi}
    \int_{-\infty}^{+\infty}\dd\omega'\,
    G_{0,jl}^{r,<}[\omega-\omega']\,
    \frac{\delta G_{0,jl}^<[\omega']}{\delta T}.
    \label{eq:wang2007_eq89_recovered}
\end{equation}
Equation~\eqref{eq:wang2007_eq89_recovered} is precisely the paper's Eq.~(89). Combining Eq.~\eqref{eq:wang2007_eq85_recovered}, Eq.~\eqref{eq:wang2007_eq86_recovered}, Eq.~\eqref{eq:wang2007_eq88_recovered}, and Eq.~\eqref{eq:wang2007_eq89_recovered} gives the full linear-response benchmark chain behind Fig.~5.

\subsubsection{Self-consistent mean-field logic}
The same derivative logic also explains why the Wang 2007 paper is a bridge toward self-consistent mean-field theory. Once the self-energy is treated as a functional of $G$, any change in temperature, bias, or displacement covariance induces
\begin{equation}
    \delta G \propto \delta\Sigma_n[G],
\end{equation}
so the problem becomes a self-consistent functional-derivative problem. This is exactly the logical bridge from the cubic LO benchmark of Fig.~5 to the later SCMF quartic closures.

\subsection{Recovery of the 2008 review and journal version (\texttt{0802.2761v1.pdf}, \texttt{transp.pdf})}
These two files are the review formulation of the ballistic theory. Their core formulas are recovered directly from the harmonic current derivation.

Start from the exact current formula in the harmonic limit,
\begin{equation}
    I_L = -\int \frac{\dd\omega}{2\pi}\,\hbar\omega\,\Tr\sqbrak{G_0^r\Sigma_L^< + G_0^<\Sigma_L^a}.
    \label{eq:ballistic_current_start_for_review}
\end{equation}
Use the harmonic lesser Green's function
\begin{equation}
    G_0^< = G_0^r(\Sigma_L^< + \Sigma_R^<)G_0^a.
    \label{eq:ballistic_lesser_for_review_recovery}
\end{equation}
Insert Eq.~\eqref{eq:ballistic_lesser_for_review_recovery} into Eq.~\eqref{eq:ballistic_current_start_for_review}:
\begin{align}
    I_L
    &= -\int \frac{\dd\omega}{2\pi}\,\hbar\omega\,
       \Tr\sqbrak{G_0^r\Sigma_L^< + G_0^r(\Sigma_L^< + \Sigma_R^<)G_0^a\Sigma_L^a}.
\end{align}
Now use
\begin{equation}
    \Sigma_\alpha^< = -\ii f_\alpha\Gamma_\alpha,
    \qquad
    \Gamma_\alpha = \ii(\Sigma_\alpha^r - \Sigma_\alpha^a),
\end{equation}
and standard trace rearrangements. After cancellation of the equilibrium pieces, the net two-terminal current becomes
\begin{equation}
    I = \int_0^{+\infty}\frac{\dd\omega}{2\pi}\,\hbar\omega\,\Tr\sqbrak{G_0^r\Gamma_L G_0^a\Gamma_R}\sqbrak{f_L-f_R}.
\end{equation}
Thus we recover the Landauer formula
\begin{equation}
    I = \int_0^{+\infty}\frac{\dd\omega}{2\pi}\,\hbar\omega\,\mathcal T(\omega)\sqbrak{f_L-f_R},
\end{equation}
with Caroli transmission
\begin{equation}
    \mathcal T(\omega) = \Tr\sqbrak{G_0^r\Gamma_L G_0^a\Gamma_R}.
    \label{eq:caroli_recovered_from_review}
\end{equation}
This is exactly the central content of the 2008 review and its journal version.

The review's emphasis on surface Green's functions is simply the computational side of the same formalism. Since
\begin{equation}
    \Sigma_\alpha^r = V^{C\alpha}g_\alpha^rV^{\alpha C},
\end{equation}
the only extra numerical ingredient needed for the ballistic problem is a stable method for evaluating the surface Green's functions $g_\alpha^r$ of the semi-infinite leads.

\subsection{Recovery of Xu \emph{et al.} 2008 (\texttt{0807.3819v1.pdf})}
Xu \emph{et al.} reformulate the same cubic interaction physics in a mode-space quasiparticle language. We recover their main formulas by starting from the general cubic self-energy and then projecting onto one mode.

From the cubic interaction section we already obtained the general structure
\begin{equation}
    \Sigma^{(3)} \propto \Phi^{(3)} G G \Phi^{(3)}.
    \label{eq:xu_starting_cubic_structure}
\end{equation}
Project Eq.~\eqref{eq:xu_starting_cubic_structure} onto a harmonic mode basis. Then the retarded self-energy of mode $\lambda$ is built from pairs of intermediate modes $\lambda_1,\lambda_2$:
\begin{equation}
    \Sigma_\lambda^{r,(3)}(\omega)
    =
    \sum_{\lambda_1\lambda_2}
    \abs{V^{(3)}_{\lambda\lambda_1\lambda_2}}^2
    \mathcal I_{\lambda_1\lambda_2}(\omega),
    \label{eq:xu_mode_retarded_self_energy_recovered}
\end{equation}
where $\mathcal I_{\lambda_1\lambda_2}(\omega)$ is the frequency integral built from the two internal propagators.

Now recover the frequency shift and linewidth. The dressed pole condition is
\begin{equation}
    \omega^2 - \omega_\lambda^2 - \Sigma_\lambda^{r,(3)}(\omega) = 0.
\end{equation}
Write the pole as
\begin{equation}
    \omega = \omega_\lambda + \Delta\omega_\lambda - \ii\Gamma_\lambda.
\end{equation}
Then
\begin{equation}
    \omega^2 \approx \omega_\lambda^2 + 2\omega_\lambda\Delta\omega_\lambda - 2\ii\omega_\lambda\Gamma_\lambda.
\end{equation}
Substitute this into the pole equation and separate real and imaginary parts:
\begin{align}
    \Delta\omega_\lambda &\approx \frac{\re\Sigma_\lambda^{r,(3)}(\omega_\lambda)}{2\omega_\lambda}, \\
    \Gamma_\lambda &\approx -\frac{\im\Sigma_\lambda^{r,(3)}(\omega_\lambda)}{2\omega_\lambda}.
    \label{eq:xu_shift_linewidth_recovered}
\end{align}
Therefore
\begin{equation}
    \tau_\lambda^{-1} = 2\Gamma_\lambda.
    \label{eq:xu_lifetime_recovered}
\end{equation}
This is the explicit recovery of the Xu lifetime formula from the retarded self-energy.

Once the lifetime is known, the mean free path is
\begin{equation}
    \ell_\lambda = v_\lambda \tau_\lambda,
    \label{eq:xu_mean_free_path_recovered}
\end{equation}
with $v_\lambda$ the group velocity. The Xu paper then uses $\ell_\lambda$ to interpolate from ballistic to diffusive transmission. The conceptual chain is therefore
\begin{equation}
    \Phi^{(3)}
    \Rightarrow
    \Sigma_\lambda^r
    \Rightarrow
    \Gamma_\lambda,\tau_\lambda,\ell_\lambda
    \Rightarrow
    \mathcal T(\omega,L).
\end{equation}
This shows that the Xu approach is not separate from the Wang approach; it is the mode-space interpretation of the same cubic LO NEGF machinery.

\subsection{Recovery of Wang 2013 review (\texttt{1303.7317v1.pdf})}
The 2013 review is the broadest conceptual summary in the local set. The key formula we need to recover is the quartic SCMF closure.

Start from the quartic Hamiltonian in Wang's notation,
\begin{equation}
    H_n^{(4)} = \frac{1}{4}\sum_{ijkl}T_{ijkl}u_i u_j u_k u_l.
    \label{eq:wang2013_quartic_start}
\end{equation}
Now apply the quartic pair decoupling. For a symmetric quartic tensor, there are three distinct pairings that contribute equally to a two-point kernel:
\begin{equation}
    (ij)(kl),
    \qquad
    (ik)(jl),
    \qquad
    (il)(jk).
\end{equation}
Therefore the induced mean-field two-point operator is proportional to
\begin{equation}
    3\sum_{kl}T_{jj'kl}\braket{u_k u_l}.
\end{equation}
Convert the equal-time correlator to a lesser Green's function using
\begin{equation}
    \braket{u_k u_l} = \ii\hbar G_{kl}(\tau,\tau)
\end{equation}
in contour notation. Since the contraction comes from one local quartic vertex, the result is local in contour time. Thus
\begin{equation}
    \Sigma_{n,jj'}(\tau,\tau') = 3\ii\hbar\,\delta(\tau,\tau')\sum_{kl}T_{jj'kl}G_{kl}(\tau,\tau).
    \label{eq:wang2013_quartic_scmf_recovered}
\end{equation}
Equation~\eqref{eq:wang2013_quartic_scmf_recovered} is exactly the quartic SCMF closure used in the review and in the Fig.~4 benchmark family.

The paper arrives at Eq.~\eqref{eq:wang2013_quartic_scmf_recovered} through the equation-of-motion hierarchy. In the review's notation, the exact two-point equation contains the four-point Green's function:
\begin{equation}
    \sqbrak{\paren{I\partial_{\tau_1}^2 + K^C + \Sigma}G(\tau_1,\tau_2)}_{j_1j_2}
    =
    -\delta(\tau_1,\tau_2)\delta_{j_1j_2}
    -
    \sum_{j_3j_4j_5}
    T_{j_1j_3j_4j_5}
    G_{j_3j_4j_5j_2}(\tau_1,\tau_1,\tau_1,\tau_2).
    \label{eq:wang2013_eq122_recovered}
\end{equation}
Equation~\eqref{eq:wang2013_eq122_recovered} is exact. The complication is that the two-point function is coupled to a four-point function, which in turn would couple to a six-point function, producing the BBGKY-like hierarchy.

The review then rewrites Eq.~\eqref{eq:wang2013_eq122_recovered} in integral form,
\begin{equation}
    G(1,2)
    =
    G_0(1,2)
    +
    \int d3\,d4\,d5\,d6\,
    G_0(1,3)\,
    T(3,4,5,6)\,
    G(4,5,6,2),
    \label{eq:wang2013_eq124_recovered}
\end{equation}
where the contour-time-dependent coupling $T(3,4,5,6)$ contains the contour delta functions needed to synchronize the four interaction times. Equation~\eqref{eq:wang2013_eq124_recovered} is simply the integral form of the exact equation of motion.

To close the hierarchy, the review approximates the four-point Green's function by a symmetrized product of two-point Green's functions, i.e. a Wick-like factorization:
\begin{equation}
    G(1,2,3,4)
    \approx
    \text{symmetrized products of } G(1,2)G(3,4),\;
    G(1,3)G(2,4),\;
    G(1,4)G(2,3).
    \label{eq:wang2013_eq125_recovered_symbolically}
\end{equation}
The exact prefactor depends on the review's Green's-function normalization, but the structural point is clear: the four-point function is closed in terms of two-point functions. Substitute Eq.~\eqref{eq:wang2013_eq125_recovered_symbolically} back into Eq.~\eqref{eq:wang2013_eq122_recovered} or Eq.~\eqref{eq:wang2013_eq124_recovered}. Because the quartic tensor is symmetric, the three possible pairings contribute equally to the two-point kernel. That is exactly where the prefactor $3$ in Eq.~\eqref{eq:wang2013_quartic_scmf_recovered} comes from. The remaining equal-time correlator becomes $i\hbar G_{kl}(\tau,\tau)$, giving the explicit self-consistent nonlinear self-energy of Eq.~\eqref{eq:wang2013_quartic_scmf_recovered}.

The review makes an important physical comment here: this nonlinear self-energy is real. Therefore it only shifts the frequencies and does not generate a finite phonon lifetime. This is precisely why the approximation is called self-consistent mean field rather than a full scattering theory.

For the Fig.~4 benchmark, the review then specifies the bath model and the one- and two-particle quartic parameters explicitly. The bath retarded self-energy is written as
\begin{equation}
    \Sigma_\alpha^r[\omega]
    =
    \frac{1}{\pi}\,P\!\!\int_{-\infty}^{+\infty}\frac{J_\alpha(\omega')}{\omega-\omega'}\,\dd\omega'
    - iJ_\alpha(\omega),
    \qquad \alpha=L,R,
    \label{eq:wang2013_lorentz_drude_self_energy}
\end{equation}
with Lorentz--Drude spectral density
\begin{equation}
    J_\alpha(\omega) = \frac{\epsilon^2\omega}{1+\omega^2/\omega_D^2}.
    \label{eq:wang2013_lorentz_drude_J}
\end{equation}
The review caption then fixes
\begin{equation}
    \epsilon = 6.0321\;{\rm meV}/(\text{\AA}^2 u),
    \qquad
    \hbar\omega_D = 10\;{\rm eV},
    \qquad
    T_L = 1.25T,
    \qquad
    T_R = 0.75T.
    \label{eq:wang2013_fig4_bath_parameters}
\end{equation}
For the one-particle benchmark, the review uses
\begin{equation}
    \Omega^2 = 60.321\;{\rm meV}/(\text{\AA}^2 u),
    \qquad
    T_{1111} = 0.241,\;1.2,\;2.4\;{\rm eV}/(\text{\AA}^4 u^2),
    \label{eq:wang2013_fig4_one_particle_params}
\end{equation}
corresponding to the black, red, and blue curves. For the two-particle benchmark it uses
\begin{equation}
    K_{11}=K_{22}=60.321,
    \qquad
    K_{12}=K_{21}=-30.165
    \qquad
    {\rm meV}/(\text{\AA}^2 u),
    \label{eq:wang2013_fig4_two_particle_K}
\end{equation}
and quartic tensors
\begin{align}
    &T_{1111}=T_{2222}=0.483,
    \qquad
    T_{\{1,1,1,2\}}=-T_{\{1,1,2,2\}}=-0.241,
    \label{eq:wang2013_fig4_two_particle_black} \\
    &T_{1111}=T_{2222}=2.4,
    \qquad
    T_{\{1,1,1,2\}}=-T_{\{1,1,2,2\}}=-1.2,
    \label{eq:wang2013_fig4_two_particle_red} \\
    &T_{1111}=T_{2222}=4.8,
    \qquad
    T_{\{1,1,1,2\}}=-T_{\{1,1,2,2\}}=-2.4,
    \label{eq:wang2013_fig4_two_particle_blue}
\end{align}
in units of ${\rm eV}/(\text{\AA}^4 u^2)$, where the curly braces indicate all permutations of the indices. These are the exact benchmark-facing parameters that should be enforced when reproducing the review's Fig.~4.

Finally, the review states that the current is still calculated with the Caroli/Landauer formula, but now with the nonlinear retarded self-energy incorporated into $G^r$. In our notation this means
\begin{equation}
    G^r_{\rm SCMF}
    =
    \sqbrak{(G_0^r)^{-1} - \Sigma_n^{r,\rm SCMF}}^{-1},
    \label{eq:wang2013_scmf_dressed_green}
\end{equation}
followed by
\begin{equation}
    \mathcal T_{\rm SCMF}(\omega)
    =
    \Tr\sqbrak{G_{\rm SCMF}^r\Gamma_L G_{\rm SCMF}^a\Gamma_R},
    \label{eq:wang2013_scmf_caroli}
\end{equation}
and then the Landauer current integral. This makes the Fig.~4 benchmark chain completely explicit:
\begin{equation}
    \text{quartic Hamiltonian}
    \Rightarrow
    \text{Eq.~\eqref{eq:wang2013_quartic_scmf_recovered}}
    \Rightarrow
    G_{\rm SCMF}^r
    \Rightarrow
    \mathcal T_{\rm SCMF}
    \Rightarrow
    I(T).
    \label{eq:wang2013_fig4_benchmark_chain}
\end{equation}

The review also discusses lowest-order cubic lifetimes. Those are recovered exactly as in the Xu section above, which confirms an important separation of approximation roles in the literature: cubic interactions are naturally introduced through LO or SCBA scattering, while quartic interactions are naturally introduced through Hartree or SCMF renormalization.

\subsection{What these papers imply for future material calculations}
Taken together, the local papers imply a clean hierarchy of reusable theory.
\begin{enumerate}[label=\arabic*.]
\item Ballistic transport is governed by the exact harmonic center Green's function, the lead self-energies, and the Landauer--Caroli current formula.
\item Cubic interactions enter naturally through the two-vertex cubic self-energy, first at LO and then at SCBA if feedback is needed.
\item Quartic interactions enter naturally through the one-vertex Hartree contraction, first at MF/SCMF and later, if needed, through more expensive dynamic quartic self-energies.
\item Paper-specific toy-model equations should be understood as special sparse realizations of the same tensor-level formalism.
\end{enumerate}
So the literature set supports one general tensor-based inelastic NEGF theory, together with several benchmark-specific wrappers, rather than several unrelated theories.

\subsection{Summary of recovered formulas}
For implementation purposes, the main recovered formulas are:
\begin{align}
    G_0^r &= \sqbrak{(\omega+\ii\eta)^2 I - K^C - \Sigma_L^r - \Sigma_R^r}^{-1}, \\
    G^r &= \sqbrak{(G_0^r)^{-1} - \Sigma_n^r}^{-1}, \\
    G^< &= G^r\paren{\Sigma_L^< + \Sigma_R^< + \Sigma_n^<}G^a, \\
    I_L &= -\int \frac{\dd\omega}{2\pi}\,\hbar\omega\,\Tr\sqbrak{G^r\Sigma_L^< + G^<\Sigma_L^a}, \\
    \mathcal T_{\rm bal}(\omega) &= \Tr\sqbrak{G_0^r\Gamma_L G_0^a\Gamma_R}, \\
    \Sigma^{(3)}_{\rm LO} &\text{: obtained from the two-vertex cubic contraction}, \\
    \Sigma^{(4)}_{\rm MF} &\text{: obtained from the one-vertex quartic Hartree contraction}, \\
    \Sigma^{\rm SCBA} &\text{: obtained by evaluating the Born diagram class with dressed } G.
\end{align}
Every central working equation in the local literature set is one of these formulas directly or a benchmark-specific rewriting of one of them.


\section{Concluding Remarks}
The central theoretical picture that emerges from the literature set is the following.

\begin{enumerate}[label=\arabic*.]
\item The harmonic phonon transport problem is completely determined by the mass-normalized dynamical matrix, lead self-energies, and the Keldysh current formula.
\item Anharmonicity enters through nonlinear self-energies. For cubic interactions the lowest nontrivial dynamic scattering appears at second order in the cubic vertex; for quartic interactions the static Hartree renormalization already appears at first order in the quartic vertex.
\item LO, MF, and SCBA are best viewed as different closures of the Dyson--Keldysh system rather than unrelated methods.
\item The Wang 2006/2007/2013 family of papers provides a continuous hierarchy: current formula, contour formalism, perturbation theory, and self-consistent mean-field closures all live inside one common framework.
\item The Xu \emph{et al.} 2008 momentum-space formulation connects the same NEGF machinery to phonon lifetime, frequency shift, and mean-free-path extraction. That link is important when bridging from toy chains to materials.
\end{enumerate}

For future code development, the clean separation should remain:
\begin{itemize}
\item general solver abstractions should be written in terms of tensors $K$, $\Phi^{(3)}$, $\Phi^{(4)}$, and lead self-energies;
\item benchmark-specific papers should live in thin wrappers that enforce the paper's units, toy-model Hamiltonian, and plotting conventions.
\end{itemize}

\end{document}
